\documentclass[aps,rmp,groupedaddress,twocolumn,10pt]{revtex4-2}

\usepackage{graphicx}
\usepackage{hyperref}
\usepackage{amssymb}
\usepackage{amsmath}
\usepackage[german,english]{babel}

\begin{document}

\title{Existence, relativity, and cosmology: reconciling foundational tensions in physics}

\author{Daryl Janzen}
\email[]{daryl.janzen@usask.ca}
\affiliation{Department of Physics and Engineering Physics, University of Saskatchewan, 116 Science Place, Saskatoon, SK S7N 5E2, Canada}

\date{\today}

\begin{abstract}
This paper addresses foundational tensions in modern physics by reexamining the concept of existence through the lenses of relativity, black hole physics, and cosmology. It integrates innovations in geometry, philosophy, and linguistics to critique entrenched biases and linguistic imprecision that have historically constrained the development of coherent physical theories. Challenging traditional arguments for the block universe, it identifies logical flaws tied to implicit meta-time assumptions and introduces formal definitions of ``existence'' consistent with a three-dimensional, enduring universe described by relativistic cosmology. A layered geometrical framework unifies cosmic time, space, and expansion, reconciling relativistic principles with empirical evidence for a cosmological rest frame and absolute cosmic time. This reinterpretation challenges Einstein's operational approach to simultaneity, offering a cosmological perspective that resolves inconsistencies between relativistic physics and cosmology. Gravitational collapse is fundamentally reexamined, revealing that fully-formed black holes with event horizons or singularities do not form within finite cosmic time, undermining assumptions behind Hawking radiation and black hole entropy. Building on these insights, the paper develops a new cosmological model that naturally resolves the expansion, horizon, flatness, dark energy, and cosmological constant problems while providing testable predictions to address the Hubble tension, early structure formation, and the origins of the universe. By embedding conservation laws and the universality of physics at its core, this framework synthesizes empirical evidence with relativistic principles, demonstrating how rethinking foundational assumptions can uncover new physics and offering a unified pathway for resolving key tensions in modern cosmology.
\end{abstract}

%\keywords{Relativity theory,Cosmic time,Block universe,Black hole physics,Cosmic microwave background (CMB),Gravitational collapse,Ontology of existence,Noumenology and phenomenology,Philosophy of science,General relativity}

\maketitle

\tableofcontents

\section{Introduction}\label{sec:intro}

\subsection{Framing the inquiry: language, bias, and foundations}

The interplay between empirical evidence, theoretical constructs, and the implicit biases of language shapes how we understand physical reality. This paper addresses the necessity for clarity in describing existence and simultaneity within physical theories. By examining the linguistic and conceptual foundations of relativity and cosmology, it explores how these frameworks both reveal and constrain our understanding of the universe.

Linguistic imprecision, particularly in the copular verb ``is,'' has led to entrenched biases and misleading concepts about the existence of physical objects and variables. Such ambiguity has hindered our understanding of time and shaped the conceptual models of reality adopted by physicists and philosophers. By refining the lexicon of existence---distinguishing when and how to describe entities as ``being'' or ``existing''---this paper proposes a coherent relativistic framework that reconciles intuition with empirical facts. This approach has the potential to resolve longstanding problems in modern cosmology, from the cosmological constant problem \cite{Weinberg1989,Carroll1992} to the Hubble tension \cite{Freedman2017,Verde2019,Freedman2023,Riess2024}, while offering a falsifiable framework for empirical testing.

This thematic framing underscores the epistemological challenges at the heart of this paper: by analysing the interplay between language, foundational assumptions, and empirical data, we address the entrenched linguistic, conceptual, and cultural biases that constrain the development of clear physical theories. As \citet{strassler2024} observed in the context of terms like ``atom'' and ``particle,'' entrenched language often limits our ability to conceptualise phenomena clearly: ``The terms that initially seem appropriate become entrenched long before they prove inadequate or misleading\ldots. Our best bet\ldots is to appreciate the dialect's limitations and find ways to circumvent them.''

Crucially, this challenge extends beyond linguistic imprecision to encompass biases embedded within intellectual paradigms across disciplines. Overcoming these biases requires revisiting foundational principles and fostering intellectual openness. This process mirrors the development of relativistic frameworks in physics, which emerged well before the implications of a four-dimensional block universe were fully realised \cite{Rietdijk1966,Putnam1967} and prior to the discovery of empirical evidence for cosmic expansion and a universal time. The historical lag underscores the interplay between empirical discoveries, mathematical formalism, and the conceptual frameworks we impose on physical reality.

\subsection{Reconciling perspectives: historical and modern challenges}

Physicists and philosophers have, for millennia, grappled with understanding the nature of time and what it means for something to physically ``exist.'' The confusion surrounding these concepts has persisted into the era of modern physics, where relativity has compounded their complexity. This paper builds on these historical efforts by addressing the tacit understanding of time that underpins much of human intuition. By reconciling this tacit knowledge with empirical data and theoretical constructs, the framework developed here offers a pathway to resolving foundational tensions in physics, integrating philosophical clarity with scientific precision to reimagine the principles of time, space, and existence.

Subjective biases further complicate the discourse, particularly in shaping interpretations of simultaneity and existence. Physicists and philosophers are human, shaped by inclinations that can influence their perspectives on reality. For instance, there is a tendency to prefer a mysterious or complex reality over a straightforward one. This predilection stems from various sources: natural scientific skepticism (``we can never truly know for certain, so it's important to explore all possible alternatives''); the allure of counterintuitive possibilities (``\ldots and wouldn't it be fascinating if a fundamental aspect of apparent reality, one we all tacitly understand, is really an illusion''); and even human ego (``\ldots and while I wouldn't argue that this is a good reason, the fact that the brightest minds have not been able to properly explicate our tacit understanding of time does make one suspicious of it''). Such biases are reinforced by an academic growth mindset, which views challenges as opportunities for intellectual advancement.

This mindset, however, can sometimes conflate depth with difficulty. As \cite{Smolin1997} observed, there was ``a peculiarity of the twentieth century that we like to believe that what is deep must also be difficult.'' He compared general relativity---which captures profound insights about our universe but is often expressed in advanced mathematics inaccessible to most undergraduates---to inaccessible literary masterpieces like \textit{Finnegan's Wake} or abstract postmodern art. This conflation may explain why interpretations that lean into the counterintuitive or the mysterious---such as the block universe---tend to gain traction.

Despite these subjective influences, a prevailing consensus has emerged. Many physicists \cite{barbour2001,davies2002,Greene2004,rovelli2019} and philosophers \cite{callender2017} agree with Einstein’s claim that ``the distinction between past, present, and future is only an illusion'' \cite{Schilpp1959}. However, this consensus rests on an ill-conceived notion of ``existence'' in modern physics, which has shaped standard interpretations of relativity, including the block universe (introduced in \S~\ref{sec:block} and discussed in \S~\ref{sec:existence}), the meaning of relativity (discussed in \S~\ref{sec:relativity}), and the nature of black holes (discussed in \S~\ref{sec:BH}). By revisiting these concepts, this paper critiques the assumptions underlying Einstein’s interpretation and its operationalist extensions, laying the groundwork for a more coherent and intuitive framework.

\subsection{Cosmic reality: evidence and interpretation}

Logical analysis of relativistic existence, accounting for foundational principles, theoretical inferences, and empirical evidence from relativity and cosmology (as discussed in \S~\ref{sec:relativity} and \S\S~\ref{sec:block}-\ref{sec:existence}), reveals a profound cognitive dissonance within modern physics. This dissonance stems from the compartmentalization of conflicting worldviews: on one hand, the operationalist interpretation of relativity; on the other, empirical evidence from cosmology that points to a preferred rest frame and an absolute cosmic time. This paper critiques how these compartmentalized perspectives obscure the deeper inconsistencies in foundational assumptions, perpetuating a fragmented understanding of existence.

Relativity is frequently interpreted as negating an objective ``now'' and asserting that simultaneity is relative. This interpretation relies on an operational definition of simultaneity, where ``simultaneous'' events are those that occur synchronously within a given reference frame. Further, it assumes that no observational evidence can identify a ``preferred'' state of rest, making the operational definition the most ``objective'' available. This perspective has deeply shaped conceptions of existence, fostering the conclusion that time is an illusion and reality is best described as a block universe. It also underpins other theoretical constructs, such as the standard interpretation of gravitational collapse (leading to black holes with event horizons and singularities thought to exist in our universe) and speculative possibilities like relativistic time travel. Over the past century, the acceptance of these ideas has created significant intellectual inertia, making it difficult to challenge the premise that a ``preferred'' state of rest cannot be observed.

Modern astronomy and cosmology present a starkly different worldview, grounded in a well-defined cosmic ``now'' and an associated ``preferred'' state of rest. Governed by general relativity, the universe described by cosmologists is a three-dimensional one, emerging from a big bang roughly 13.8 billion years ago and expanding ever since. Over the past century---and especially in the twenty-first century---our understanding of this universe has become remarkably precise. Observations of the cosmic microwave background (CMB), the oldest light emitted when the universe was 380,000 years old, integrate with models of expansion, structure formation, and light-matter interactions during various epochs. These models describe key transitions, supported by evidence: from an early plasma state to recombination, when neutral atoms formed and photons began traveling freely; through reionisation, driven by early stars and galaxies ionising the intergalactic medium; and from an initially decelerating expansion to the current era of acceleration, driven by dark energy, which began roughly around the Solar System's formation.

These two perspectives, with their contrasting worldviews, are deeply intertwined. Both propose significant implications for the nature of space, time, and physical reality. While the standard cosmological model describes a cosmic time that is absolute, with a preferred ``now'' and associated cosmic ``simultaneity,'' this interpretation conflicts with relativity's operational definition. Textbooks dedicated to cosmology typically begin with treatments of relativity, framing it as the fundamental theory underpinning the standard cosmological model \cite{Ryden2017}. This hierarchical framing can give the impression that the preferred frame of cosmology holds only practical significance, useful for describing large-scale phenomena but lacking fundamental meaning.

From an objective standpoint, however, this is not the logical structure of cosmological science. The standard cosmological model begins with the principle that there is a cosmic time. This assumption is essential to define ``space'' as expanding, which is the standard explanation of the observed redshift-distance relation. Anticipating that the model must be relativistic, the operational definition of simultaneity is assumed valid in the cosmic rest frame, leading to the further assumption that space is synchronous in this frame. Additional assumptions of isotropy and homogeneity constrain the spatial slices to be maximally symmetric, with positive, negative, or zero curvature. Finally, ensuring that this space-time is a solution to Einstein's field equations constrains the form of cosmic expansion, linking it to the spatial curvature, energy densities, pressures, and current expansion rate. These constraints yield the Friedmann-Lemaître-Robertson-Walker (FLRW) cosmological models \cite{Friedmann1922,Friedmann1924,Lemaitre1927,Lemaitre1931a,Lemaitre1931b,Lemaitre1931c,Robertson1936a,Robertson1936b,Walker1935a,Walker1935b}, which form the basic geometrodynamical framework of modern cosmology.

The success of these models in explaining a wide variety of phenomena---from the detailed anisotropy of the CMB \cite{Aghanim2020} to the local expansion rate of space \cite{Freedman2017,Freedman2023,Riess2024} and the formation and evolution of galaxies and clusters \cite{Springel2005,Springel2006,Haslbauer2022}---provides strong evidence for their validity. Crucially, the evidence supporting the hypothesis of an expanding universe also supports the more fundamental hypothesis that a cosmic time exists through which space is defined and its expansion is described. By the basic principles of scientific reasoning, it is logically inconsistent to accept the empirical evidence for the expansion of the universe while rejecting the more fundamental hypothesis that underpins it: the existence of cosmic time. Furthermore, the standard interpretation of the CMB dipole anisotropy is that it provides an accurate measurement of both our speed and direction of motion through the universe ($v=369.82\pm0.11$~km~s$^{-1}$ towards the constellation Crater \cite{Aghanim2020a}). Thus, it is equally inconsistent to maintain that our velocity relative to a ``preferred'' rest frame cannot be observed.

From the perspective of modern cosmology, it seems that rather than being a convenient but secondary construct, the coordinate frame of the standard cosmological model has objective significance. The scientific evidence suggests our universe does possess a well-defined ``now,'' even though events at the same cosmic time are not generally synchronous in local reference frames. This indicates that the operational definition of simultaneity should not be granted ontological primacy and that relative synchronicity lacks fundamental significance in describing physical ``existence.''

Reconciling these conflicting worldviews leads to the conclusion that the enduring three-dimensional universe described by astronomers and cosmologists provides a more accurate representation of physical reality than the four-dimensional existence often associated with relativity, in which space-time is thought of as the basic fabric of reality. This reinterpretation suggests that relativity itself should be fundamentally understood cosmologically, as providing the metrical description of events within a well-defined, enduring three-dimensional space---our universe---in both local and cosmological frames.

\subsection{Toward a unified framework}

Such a significant reinterpretation of one of our foundational theories of physical reality warrants careful scrutiny to ensure the associated concepts and definitions are as clear and accurate as possible. This paper therefore presents a fresh analysis, beginning in \S\S~\ref{sec:block}-\ref{sec:relativity} by developing a concept of relativistic existence, together with explicit language to describe it, that aligns with the empirical evidence and our basic intuition about the natures of time and space. 

Since the \textit{learned} intuition of modern physicists tends to be influenced by notions of space-time's existence, possibly as a block universe, I start in \S~\ref{sec:block} by outlining this counter-position through three arguments that broadly shaped this perspective. In \S~\ref{sec:existence}, I address the common problem with such arguments: a misunderstanding of the physical meaning of the word ``exist.'' I propose a stripped-down physical definition of existence, independent of life, consciousness, or thermodynamics, to clarify this concept. Alongside this definition, I introduce related terms such as ``happen,'' ``event,'' ``moment,'' and ``space-time,'' carefully delineating their roles in physical contexts. Consistent with \citeauthor{strassler2024}'s~\citeyearpar{strassler2024} recommendation to circumvent linguistic limitations in physics, I advocate for the use of new, non-existential copulas to describe non-existential objects (e.g., space-time events or the block universe). \S~\ref{sec:existence} concludes with a discussion of how each of the examples from \S~\ref{sec:block} misuses the concept of existence, fostering misunderstanding both of what it means to exist physically and of the theories they seek to support.

In \S\S~\ref{sec:relativity} and \ref{sec:BH}, I address two significant objections to the explicit concepts of time and existence developed in \S~\ref{sec:existence}: first, that they conflict with standard descriptions of the relativity of simultaneity; and second, that they are inconsistent with a universe containing black holes. In \S~\ref{sec:relativity}, I re-examine the original motivation for Einstein's interpretation of relativity and argue that it no longer holds up against current evidence. Instead, relativity should be fundamentally interpreted cosmologically. In \S~\ref{sec:BH}, I find that the standard model of gravitational collapse to a black hole cannot be justified within the cosmological interpretation of relativity. Gravitationally collapsing matter, instead, is better described as asymptotically approaching its event horizon.

The cosmological interpretation of relativity developed in \S\S~\ref{sec:block}-\ref{sec:BH} offers a fundamentally different picture of physical reality. In this view, space-time is ultimately not ``reality'' at all; rather, it is the four-dimensional map of events that unfold in a three-dimensional reality that \textit{exists}---such as the one modern cosmologists describe, which began with a big bang and is now 13.8 billion years old. According to this interpretation, the three-dimensional universe's existence metrically unfolds in accordance with the principles of general relativity. 

This interpretation differs fundamentally from Einstein's original conception of relativity and suggests empirical differences. Already, we've noted that the cosmological interpretation is supported by evidence falsifying a key assertion of operational relativity: that motion relative to a preferred cosmic rest frame could never be observed. Yet modern cosmology faces several challenges, with the number and significance of known problems actually increasing in the era of precision measurements. Long-standing problems include the lack of explanation for \textit{why} the universe expands \cite{Eddington1933},\footnote{Specifically, as \citet{Eddington1933} framed the problem in a comment on the model of \citet{Einstein1932} (which incidentally became accepted as the standard cosmological model throughout much of the 20th century): ``There are only two ways of accounting for large receding velocities of the nebulae: (1) They have been produced by an outward directed force as we here suppose, or (2) as large or larger velocities have existed from the beginning of the present order of things. Several rival explanations of the recession of the nebulae, which do not accept it as evidence of the repulsive force, have been put forward. These necessarily adopt the second alternative, and postulate that the large velocities have existed from the beginning. This might be true; but can scarcely be called an \textit{explanation} of the large velocities.\ldots The theory recently suggested by Einstein and de Sitter, that in the beginning all the matter created was projected with a radial motion so as to disperse even faster than the present rate of dispersal of galaxies (they do not state this in words, but it is the meaning of their mathematical formulae), leaves me cold. One cannot deny the possibility, but it is difficult to see what mental satisfaction such a theory is supposed to afford.''} the dominance of invisible components like dark matter and dark energy in the measured rate of cosmic expansion \cite{Riess1998,Perlmutter1999}, and the unresolved cosmological constant problem, where vacuum energy predictions exceed observed values by 120 orders of magnitude \cite{Weinberg1989,Carroll1992}.

Moreover, tensions have emerged between early- and late-universe measurements of the Hubble constant $H_0$ \cite{Freedman2017,Freedman2023,Riess2024}, alongside difficulties in explaining significant structures observed in the early universe with current models of structure formation \cite{Haslbauer2022}. Inflation, the prevailing explanation for the horizon and flatness problems, faces its own explanatory challenges, as the conditions for producing a universe with observed properties appear unnatural within a generic inflationary framework \cite{Ijjas2014,Ijjas2017,Garfinkle2023}.

In the pursuit of new physics, the path forward may not lie in tweaking established paradigms or introducing ad hoc fields like inflation or early dark energy \cite{Kamionkowski2023} to address these tensions. While such approaches have yielded provisional solutions, they remain constrained within the same foundational framework that has repeatedly shown its limitations. Instead, these challenges may indicate a deeper issue, suggesting that the next breakthrough in physics requires a rethinking of foundational principles.

In particular, \S~\ref{sec:cosmology} introduces a new cosmological model grounded in geometrical principles informed by the cosmological interpretation of relativity. By rethinking foundational assumptions in FLRW cosmology, the model resolves persistent tensions that have long challenged the standard framework. These include the flatness, horizon, and cosmological constant problems, which are shown to emerge naturally from the model's treatment of cosmic time, expansion, and structure as interconnected within a single geometric framework. In contrast to ad hoc mechanisms like inflation or early dark energy, this approach provides a coherent explanation grounded in first principles. The model also offers testable predictions, such as recalibrated density parameters and a reinterpretation of the sound horizon, which provide new avenues for resolving discrepancies like the Hubble tension and challenges in early structure formation.

The implications of this framework extend beyond cosmology, revealing potential connections to fundamental symmetries in particle physics. The model suggests that the duality of matter and antimatter, as well as key symmetries underlying electromagnetism and the weak interaction, may arise from the universe's geometric structure. These insights hint at a deeper unity between the forces governing the cosmos, bridging phenomena like gravitational collapse and the big bang with the emergent properties of particle interactions. This integrated perspective highlights the model's capacity to unify cosmology and particle physics within a cohesive theoretical framework, offering a path forward for addressing unresolved questions at the intersection of these fields.

These efforts culminate in five key contributions that address foundational challenges in modern physics: 
\begin{itemize}
    \item \textbf{A new understanding of physical existence:} The paper demonstrates that modern physics and philosophy have operated for over a century with a fundamentally flawed understanding of physical existence. This misunderstanding underpins key conceptual missteps in contemporary theories, revealing a critical need for reassessment.
    \item \textbf{A reinterpretation of relativity:} By challenging Einstein’s interpretation of relativity, the paper argues that modern empirical evidence supports a cosmological reinterpretation. This approach redefines simultaneity and cosmic time, addressing longstanding inconsistencies in our understanding of relativistic phenomena.
    \item \textbf{A critique of gravitational collapse:} The paper exposes a flawed concept of gravitational collapse that has persisted in black hole physics and quantum gravity research, rooted in Einstein’s misinterpretation of relativity. It provides a basis for rethinking these areas to address their unresolved challenges.
    \item \textbf{A paradigm shift in cosmology:} The paper critiques the standard cosmological model, demonstrating that it bears a fundamental flaw that parallels the basis of geocentrism, similarly leading to a reliance on ad hoc mechanisms like inflation and dark energy to reconcile foundational inconsistencies. It proposes a layered geometry hypothesis, informed by the cosmological reinterpretation of relativity, as a promising alternative framework to be explored collaboratively.
    \item \textbf{A call for epistemological reform in physics:} Identifying the prioritization of operational convenience and ad hoc hypotheses over reductive analysis as a central issue, the paper calls for an epistemological shift. It advocates a return to foundational inquiry and conceptual clarity as essential steps toward overcoming entrenched limitations and unlocking the next wave of progress in physics.
\end{itemize}

\section{The block universe}\label{sec:block}

\subsection{The concept in Ancient philosophy}

Parmenides of Elea was likely the first to form an argument in favour of what is now referred to as a block universe, or block space-time, description of physical reality. Parmenides deduced a paradoxical description of the world roughly as follows \cite{Gallow1984}. He began with the basic principle that ``what-is, is, and what-is-not, isn't.'' He noted that the present is certainly part of what-is, but also, since we can recall the past and are certain that it really happened, that the past is also part of what-is. And then, making the natural assumption that experience confirms to be true every day, that the future is also real, he noted that the future is also part of what-is. Thus, he concluded that all of eternity, from the distant past to the distant future, must all be real, as it is all part of what-is. Therefore, he found that apparent change is an illusion.

Parmenides' student, Zeno, presented several paradoxes in support of Parmenides' position that all of eternity is one and unchanging. For example, one paradox noted that motion cannot occur since in order to travel any distance one must first reach the half-way point. But then, before reaching the halfway point, one must reach one quarter of the distance, etc. Therefore, one must traverse an infinite sequence of steps, beginning with an initial, infinitesimally small step, in order to traverse any distance, which is impossible, so motion must be impossible.

\subsection{The concept in popular culture}

Modern arguments have been similarly established. For example, in the opening pages of \textit{The Time Machine}, \citet{Wells1895} drew a similar conclusion. The novel opens with the Time Traveller protagonist explaining his theory of four-dimensionalism to a group of guests. Since this was the first time travel story, Wells went into some detail, providing examples such as a set of portraits of a man throughout his life, as cross-sectional representations of his four-dimensional being. He further likened time to gravitation, noting that while we can walk around the surface of the Earth we struggle much more to move vertically through Earth's gravitational field, suggesting that a time machine may do for the time-dimension what a balloon does in transporting people into the sky.

Wells thus established a concept that has since become familiar in science fiction and fantasy literature and entertainment: the notion that the past, present and future all exist and that time travel can potentially happen. In a preface he added later, \citet{Wells1931} noted:
\begin{quote}
    That one idea is now everybody's idea. It was never the writer's own peculiar idea. Other people were coming to it. It was begotten in the writer's mind by students' discussions in the laboratories and debating society of the Royal College of Science in the [eighteen-]eighties and already it had been tried over in various forms by him before he made this particular application of it. It is the idea that Time is a fourth dimension and that the normal present is a three-dimensional section of a four-dimensional universe. The only difference between the time dimension and the others, from this point of view, lay in the movement of consciousness along it, whereby the progress of the present was constituted. Obviously there might be various ``presents'' according to the direction in which the advancing section was cut, a method of stating the conception of relativity that did not come into scientific use until a considerable time later\ldots
\end{quote}

Since this story's original publication in 1895, the concept of time travel has become a common theme, with stories tending to fall into one of two categories: either a deterministic world view, in which the events of space-time are fixed in spite of every effort made to alter them; and one in which time travellers are capable of altering events, in which case versions of the grandfather paradox are often explored. But this excerpt from Wells' later preface shows that the concept of a four-dimensional universe was not his alone, but a burgeoning one in academic circles that people were keen to accept, and it eventually became ``everybody's idea.'' He was simply the first to develop the concept into a vivid and widely read story.

\subsection{The concept as a consequence of relativity}

\textit{The Time Machine} was published the same year as Lorentz's theory of optical and electrical phenomena in moving bodies, a decade before Einstein published his first paper on the special theory of relativity \cite{Einstein1952}. A few years after that, Minkowski expounded his geometric theory of the absolute four-dimensional world, about which he said \cite{Einstein1952} ``Henceforth space by itself and time by itself will recede completely to become mere shadows and only a type of union of the two will still stand independently on its own.''

\citet{Einstein1938} described the space-time concept in a manner very similar to the description Wells had given in the opening pages of \textit{The Time Machine}. They compared a ``dynamic'' picture with a particle moving through one-dimensional space, its position changing with time, to a ``static'' one in which the particle's position is represented as a curve in a two-dimensional space-time continuum. In the static picture, they note that ``motion is represented as something which \textit{is}, which exists in the two-dimensional time-space continuum, and not as something which changes in the one-dimensional space continuum.'' They then noted that both pictures were equivalent from the perspective of classical physics, though the dynamical picture had been preferred; but that relativity changed this view, as it showed that the static picture was more convenient and provided a more objective representation of reality. They explained that the reason for this was the relativity of simultaneity, and eventually concluded: ``Even in the relativity theory we can still use the dynamic picture if we prefer it. But we must remember that this division into time and space has no objective meaning since time is no longer `absolute.' We shall still use the `dynamic' and not the `static' language in the following pages, bearing in mind its limitations.''

In fact, this final point about allowing the misuse of inappropriate language when developing the concept of a block universe will be important in our discussion in \S~\ref{sec:existence}. In terms of the current discussion, it is interesting to note that \cite{Popper2002} wrote that he once discussed indeterminism at length with Einstein, and ``tried to persuade him to give up his determinism, which amounted to the view that the world was a four-dimensional block universe in which change was a human illusion, or very nearly so. (He agreed that this had been his view, and while discussing it I called him `Parmenides'.)''

Much has been said about the connection between relativity and the concept of a block universe over the past century. As \citet{Einstein1938} noted, the main argument deduces the block universe as a consequence of the relativity of simultaneity. To summarise, imagine you and a friend are at either end of a train car moving with constant speed, and a light bulb in the centre of the car is turned on. From your perspective in the car, particularly if all the blinds are drawn on the windows so you imagine yourselves to be at rest, since the light bulb is equal distances from you and your friend, and light moves at the same finite speed in any direction, you conclude that the photons emitted when the bulb turns on reach you at the back of the car and your friend at the front of the car at \textit{the same time}, a finite interval after the bulb was turned on. Now, special relativity is based on two postulates: i.~that the speed of light (in an ideal vacuum) is the same in every frame of reference, and ii.~that the physical principles that describe the world are independent of relative inertia, and all inertial observers give equally valid descriptions of phenomena. In particular, now consider the same set of events from the perspective of an observer who stands outside the train, now with the blinds up: from this perspective, after the light turns on and the photons begin their journey towards the back and front of the car, the car also continues moving forward so that you, at the back of the car, will see that the light turned on at an earlier time since the car will have moved you slightly forward to intercept it on its way to your eye. Similarly, by the time your friend at the front of the car sees the light bulb turn on, they will have moved away from the light bulb's location at the moment it was turned on, so this event must happen at a later time. In conclusion, whether two events at separate locations are described as happening ``at the same time'' or ``one before the other'' depends on relative motion and the frame of reference in which the events are being described. This is known as the relativity of simultaneity.

Now, the deduction of a block universe follows by considering a particular moment in the frame of reference of any particular observer. From their perspective, at each point in space there may be infinitely many distant observers moving at random relative speeds, and from \textit{their} perspectives at that instant all the events ``simultaneously happening now'' extend into the past and future of the original observer. But then, given that there may be infinitely many of these other observers located all throughout space ``now,'' all moving at speeds up to the speed of light, their collective ``nows'' that are all happening ``right now'' from the initial observer's perspective cover that observer's entire past and future. Thus, all of eternity can be said to be happening ``now'' from the perspective of any observer, when considering the simultaneous perspectives of all possible other observers ``now''.

As \citet{Greene2004} put it, 
\begin{quote}
    if you buy the notion that reality consists of the things in your freeze-frame mental image right now, and if you agree that your now is no more valid than the now of someone located far away in space who can move freely, then reality encompasses all of the events in spacetime\ldots as Einstein once said, `For we convinced physicists, the distinction between past, present, and future is only an illusion, however persistent.' The only thing that's real is the whole of spacetime.

    In this way of thinking, events, regardless of when they happen from a particular perspective, just are. They all exist. They eternally occupy their particular point in spacetime. There is no flow. If you were having a great time at the stroke of midnight on New Year's Eve, 1999, you still are, since that is just one immutable location in spacetime. It is tough to accept this description, since our worldview so forcefully distinguishes between past, present, and future. But if we stare intently at this familiar temporal scheme and confront it with the cold hard facts of modern physics, its only place of refuge seems to lie within the human mind.
\end{quote}

Much like \citet{Einstein1938} had done previously, \citet{Greene2004} eventually made a note about his use of language in describing the block universe e.g.\ as ``existing'', ``eternally'', etc. As he put it, ``the feeling that time flows is deeply ingrained in our experience and thoroughly pervades our thinking and language. So much so, that we have lapsed, and will continue to lapse, into habitual, colloquial descriptions that refer to a flowing time. But don't confuse language with reality. Human language is far better at capturing human experience than at expressing deep physical laws.''

This linguistic limitation prompts us to question not only how we describe existence but also how we distinguish, both conceptually and in our formal mathematical descriptions, between what ``exists'' and what does not. The tension between linguistic precision and conceptual clarity, evident in both ancient and modern discussions, underscores the enduring challenge of accurately representing physical theories. Parmenides' assertion that ``what-is, is'' and Einstein's ``illusion of time'' may seem distant in origin, yet both grapple with the fundamental question of what it means for something to exist---a question that remains as relevant in today's physics as it was in ancient philosophy.

If we are \textit{incapable} of accurately communicating what our theories essentially describe, and to such an extent that we regularly use language we \textit{know} to be false and misleading, misrepresenting their physical meaning as these renowned physicists and physics communicators have admitted to doing, is it fair to claim that our concepts and interpretations of our theories accurately align with the formal mathematical descriptions? Of course not. The stakes are high: if linguistic imprecision distorts our conceptual frameworks, it risks misaligning theoretical understanding with empirical reality. Thus, addressing these mismatches is not a mere academic exercise but a necessary step toward ensuring that our physical theories not only describe but also reflect reality accurately. We can't both admit that our thinking about our theories is so imprecise that when we try to articulate our thoughts we knowingly use incorrect, misleading language, yet still claim that our basic concepts and our theories are correct. And so, we shall now turn our attention to developing not only a formal explicit concept of physical \textit{existence}, but also the appropriate language to describe both existential and non-existential things.

\section{Developing a formal concept of physical \textit{existence}}\label{sec:existence}

It is difficult for us to imagine, but people have not always had a clear understanding of the colour blue. In \textit{The Odyssey}, Homer famously described the sea as ``wine-dark'' rather than blue. Both the \textit{Bible} and the \textit{Rig Veda} lack direct reference to blue. Terms for ``blue'' often didn't exist in ancient languages, and blue was one of the last colours to be named as languages evolved.

The story of how the colour blue was conceptualised in ancient cultures is an interesting blend of both linguistic and perceptual history. While a clear sky is indeed ``blue,'' the colour is otherwise relatively rare to find in nature, and blue dyes were historically difficult and expensive to produce. This likely contributed to the absence of a term for the colour in early languages. While people could always perceive blue, early studies in the 19th century suggested they did not linguistically or conceptually distinguish it from other colours. The concept of blue as a separate colour category emerged as languages and cultures evolved, influenced by discoveries like the Egyptians' creation of synthetic blue pigment.

The concept of physical \textit{existence} has similarly remained unclear, even in formal academic research, and in this case under-developed language also seems to have played a significant role in our muddled understanding of what it means for something to ``exist''. As I'll argue, each of the arguments in the previous section is predicated on a critical misunderstanding, which is the basis on which they deduce the consequence of a block universe---which they then proceed to describe inaccurately as they advocate for it as the most objective consequence of theory and logic.

\subsection{Exploring concepts of existence through familiar ideas}

When we consider the concept of existence, we often draw from a mix of everyday experiences, scientific understanding, and cultural influences. Popular media---from science fiction movies to time travel tales---shapes how we think about time, space, and existence, embedding ideas that we rarely pause to examine. These familiar narratives, coupled with modern scientific paradigms, often create a patchwork of assumptions that colour our understanding of what it means for something to ``exist.'' In this subsection, I will weave together a range of thought experiments and cultural references to illuminate the various dimensions of existence, encouraging critical engagement with the concepts before delving into their formal definitions.

Imagine that Mars has suddenly blinked out of existence. It formed 4.5 billion years ago, once flowed with water, and recently became the target of robotic exploration. Then, without warning, it vanished. Could we send another rover to explore its surface? Obviously not. A place must \textit{exist} for us to travel to it. 

But what if we invented a time machine? Could we then explore Mars before it disappeared? This scenario hinges on whether Mars's 4.5 billion years of existence continues to ``exist'' in some form beyond the present moment. If not, then no time machine could reach it; Mars would simply be gone. This highlights the tension inherent in block universe arguments, which often rely on the assumption that past states like Mars's prior existence persist in some tangible form, a sort of meta-existence.

This example introduces a deeper conflict between the intuitive understanding of existence and the philosophical implications of the block universe concept. To explore this further, consider a closed FLRW universe: a 3-sphere expanding from a singularity. Cosmologists often describe this universe as existing only at the present moment, with time marking its growth. Yet proponents of the block universe argue that all its instants exist simultaneously in a four-dimensional space-time. Imagine this space-time filled with all of eternity: every three-dimensional moment of the universe's existence, from the big bang to its current state, ``existing'' together in a timeless tableau. In this mental image, a spotlight might be thought to represent the ``now,'' illuminating different slices as it sweeps across eternity \cite{Price2011}. If Mars were to vanish as suggested above, it would no longer exist in the spotlight---but it would still ``exist'' in the past, reachable by escaping the spotlight's reach. 

Now, consider a purely three-dimensional universe. Imagine traveling from one city to another. Your journey does not simply unfold across space; a map showing the route you take does not capture your entire world-tube, which unfolds over time, with each step requiring space to \textit{exist} and time to pass. This intuitive notion of existence---that entities endure and occupy space over time---forms the bedrock of our understanding of reality. If space did not exist, movement through it would be impossible. Our world's \textit{existence} is why we can move through it, interact with objects, and track the passage of time.

When we describe something as ``existing'', we confer on it the significant property that it \textit{endures} through time. A thing need not move or change in order to exist; it simply needs to \textit{be} part of our \textit{existing} three-dimensional world. A book, \textit{sitting} on a shelf, \textit{exists}, and its existence forms a four-dimensional world tube within four-dimensional space-time. You, wherever you are, contemplating what I've written in this paper: that contemplation is taking place over the course of time, as your world tube evolves.

This sense of ``existence'' is so fundamental to our understanding of reality, it is at the heart of \textit{every verb} in our language: to sit, to contemplate, to run, to fly, to \textit{be}---each of these words only makes sense in the context that the thing doing them \textit{exists}. 

Thus, when block universe proponents describe all of eternity as ``existing'' they unconsciously embed within that conception an extra dimension of time---a \textit{meta}-time dimension. This tacit dependence on meta-time contradicts the purported timelessness of the block universe, revealing a conceptual inconsistency in its foundational premise. In this view, the ``spotlight'' described above \textit{moves} through the four-dimensional block universe over the course of meta-time. And if the spotlight burned out, and the whole block just \textit{existed}, unilluminated, it would still exist in meta-time. Without meta-time, the block universe cannot ``exist'' in the enduring sense; it can only ``happen,'' with all events occurring simultaneously without any temporal endurance. This reliance on an implicit meta-time dimension underscores a critical flaw in block universe arguments: the concept's very coherence depends on assumptions that its proponents often fail to acknowledge explicitly.

This distinction between ``existing'' and ``happening'' becomes critical when we examine block universe arguments. Proponents often describe all of eternity as ``existing,'' implicitly invoking a meta-time dimension to sustain the endurance they ascribe to space-time. This concept is incorrect, as the entire block universe could only ``happen'' all-at-once. This fundamental distinction will guide our analysis of the linguistic and conceptual challenges surrounding existence.

These challenges---rooted in how we conceive of and articulate existence---underscore the importance of linguistic precision. Without a clear vocabulary to distinguish between enduring existence and transient happening, our interpretations risk perpetuating conceptual errors. To address this, we must refine the language used to describe existence and its interplay with space-time, ensuring our words accurately capture the nuances of physical reality. To reconcile these intuitive and philosophical tensions, we must turn to the precise language needed to formalise our understanding of existence in the context of modern physics.

\subsection{A linguistic framework to describe existence}

\textbf{Time} is a physical property of our three-dimensional \textbf{universe}. The universe \textbf{endures} or \textbf{exists}, and has done since the big bang, and continues to do so as it expands, which it is presently doing at an accelerating rate. 

Every \textbf{verb} describes an act that takes place over some duration because the thing doing the act \textit{exists}. For example, a person ``sits'' at a computer, or ``drives'' down the road, or ``listens'' to music, all over the course of some length of time. Thus, \textit{verbs} only apply to \textit{existing} things. Conversely, it is incorrect to say, for instance, that \cite{Einstein1938} ``motion is represented as something which \textit{is}, which exists in the two-dimensional time-space continuum'' because the time-space continuum and everything represented in it \textit{does not} exist.

An \textbf{event} iz\footnote{This is not a typo; I'll come to a definition of ``iz'' below. For now, just read it as ``is,'' although that's technically incorrect since ``is'' is a verb.} a location within our existing three-dimensional universe, at a particular instant of time. We say that events \textbf{happen} or \textbf{occur} through the course of the universe's existence Events \textit{do not} themselves exist.

\textbf{Space-time} iz the four-dimensional set of events that \textbf{happen} throughout the three-dimensional \textbf{universe} as it \textbf{exists}. Space-time \textit{does not} exist.

Existence is a non-trivial physical property of our universe and everything in it. In fact, according to relativity theory, it is a well-ordered \textit{metrical} property. As things exist, time passes in a metrically significant way, everywhere describable as events unfold in space-time. 

It is incorrect to say that space-time itself exists, as existence implies a passage of time, and this confers a meta-time within such descriptions of physical reality. Events occur over the course of time in our three-dimensional universe. They happen where and when they do, and the whole set of events that ever occur, unfolding over the course of time, forms all of space-time. To describe four-dimensional space-time as existing in the same sense that the three-dimensional universe and three-dimensional objects within it exist, is to confer upon space-time significant existential value. Just as the existence of our universe is a well-defined metrical property, so that all the events that occur in it map to a well-defined four-dimensional metric space, so should space-time's supposed existence necessitate a fifth, timelike dimension along which ``frozen'' events of space-time could be said to remain eternally unchanged.

In order to discuss things like events, instantaneous objects, and space-time, we really could do with some additional language. Consider, for example, the ``instantaneous cube'' discussed by the Time Traveller at the beginning of \textit{The Time Machine}. It is problematic to say ``the instantaneous cube is three-dimensional'' because the copula\footnote{In languages like English, we use copulas like ``is'', ``are,'' etc to couple the subject and predicate in a sentence. For example, in this sentence, the subject ``It'' is coupled to the predicate ``problematic to say\ldots'' via the copula ``is''.} ``is'' is a verb, and therefore confers existential meaning on the cube. If the cube exists, then its existence takes place over time and the events that occur throughout the three dimensions of the cube at each moment of its existence form a four-dimensional world-tube. Similarly, to say ``the block universe is four-dimensional'' conveys a non-trivial sense of existence in which all the events of the block are understood to be frozen in place, continually happening over the course of a non-trivial meta-time dimension that is distinct from the time-dimension within space-time. These linguistic imprecisions have significant consequences, shaping not only theoretical interpretations but also our broader conceptual models of reality. They underscore the need for precise language when discussing fundamental aspects of physics to avoid inadvertently conferring existential meaning on concepts that do not endure.

We encounter a similar linguistic issue when discussing the relativity of simultaneity, e.g. according to which we might say of two relatively moving observers, A and B, that ``the events happening `now' in B's reference frame are events that happened in A's past or will happen in A's future.'' In this sentence, ``are'' similarly confers a sense of existence on the set of events happening `now' from B's perspective. In this case, even an attempt to clarify by saying ``these events don't exist, they're all happening `now'{''} runs into the same problem since, in spite of the use of `now,' the statement that the events \textit{are} all happening still confers existential meaning.

These examples suggest that our language is insufficient to accurately represent the distinction between events that happen and objects that exist, since in English the copulas are verbs---the \textit{existential} verb, in fact---that explicitly confer existential meaning on the subject of a sentence.\footnote{Finally, we've dug our way down to the heart of the paper; the reason for its title choice. From here, we'll start to work our way out to form a consistent relativistic description that explicitly incorporates this deeply ingrained sense of existence.} In everyday use, this is rarely a problem since things in the world around us \textit{do} exist. But when discussing the nature of time and existence, and events that happen in the course of existence, and \textit{when formulating mathematically precise physical theories and interpreting significant implications about the nature of physical reality}, we need to be precise about whether or not the subjects of our discussion exist. Therefore, we define two new words, \textbf{iz} and \textbf{zare}, to be used in sentences involving singular and plural subjects that \textit{explicitly} do not \textit{exist}.

Having defined these non-existential copulas, the above phrases ``the instantaneous cube iz three-dimensional,'' ``the block universe iz four-dimensional,'' ``the events happening `now' in B's reference frame zare events that happened in A's past or will happen in A's future,'' and the clarifying statement, ``these events don't exist, they zare all happening `now''' are now accurate---though it still remains to provide a clear definition of ``now'' in the last statement. 

Intuitively, \textbf{now} (alternatively called the \textbf{present}) signifies the time when the word ``now'' (or ``present'') \textit{is} used---where the existential meaning of ``is'' in the definition is now deliberate. This definition captures the fact that ``now'' signifies the current time of our evolving, \textit{existing}, enduring universe, that it is a precise time, but that it is also difficult to pin down due to its fleeting nature. To illustrate, we might say ``we can never see the Sun now; we only ever see it eight minutes from now,'' meaning that at any particular time, since the Sun is eight light minutes away from Earth, it will take eight minutes for the photons emitted now by the Sun to reach Earth. In this case, since ``now'' is constantly evolving the sentence naturally updates itself.

In the next two sections, I will explain how the terminology defined above may be reconciled with relativity and, in particular, with relativistic gravitational collapse. In \S~\ref{sec:relativity}, I'll argue that Einstein's preferred interpretation of the meaning of relativity is wrong, and that a cosmological interpretation of relativity should be adopted instead, as it is more consistent with current empirical evidence and with our explicit definitions pertaining to existence (as well as with our tacit knowledge of existence, so we should not breed further misunderstanding through misuse of language). And in \S~\ref{sec:BH}, anticipating a potential objection to the cosmological interpretation of relativity if singularities were to exist, I will argue that the standard conceptual model of black hole formation from dynamic gravitational collapse is wrong, that it is the result of mistaken considerations about gravitational collapse within Einstein's interpretation of relativity, and that a cosmologically relativistic interpretation of the collapse scenario should be adopted instead.

But before getting to that, it is useful to tie up the argument in this section by explaining the specific linguistic/conceptual errors committed in each of the block universe arguments from \S~\ref{sec:block}.

\subsection{Addressing the fallacies that lead to the block universe conclusion}

In the case of the Eleatics, the fallacy in Parmenides' argument is his assumption that the past and future exist. Sure, the present universe contains evidence of the past. If you presently have a black eye, it's because something slammed hard into your face sometime in the recent past. In the present, we can also recall things about our own pasts or anticipate things that may happen in the future; but the world's record of the past and our memories and anticipations about the past and future, all exist with us in the present. The past and future do not themselves exist outside the present; outside of now. Any prior assumption about a more substantive past or future that is then used to deduce a consequence that all of eternity exists, is a circular argument; a tautology. Thus, by starting with an assumption that all the past and future \textit{exist}---that the past is part of what-is and the future is part of what-is---Parmenides came to find that all of eternity exists, which is really just what he assumed to begin with.

Regarding Zeno's paradoxes, the fallacy was that he removed the temporal value from each step and considered only sequential positions in space. In order for something to move (a verb) its spatial position must change over the course of time. Zeno was concerned with demonstrating that the number of steps required to go anywhere must be infinite, and that they must start out, or eventually become, infinitesimal, so that neither could one ever start moving (because the first step in the sequence is infinitely small) nor could they ever reach a destination (because that requires an infinite number of steps to be taken in the sequence); so the idea was that what we perceive as motion is actually a world-line that ``eternally exists'' from past to future. But for a body in motion, each of those spatial displacements has an associated temporal increment. For example, in the case of motion with constant speed, the ratio of spatial displacement to temporal increment has a constant value---and while the displacements do indeed become infinitesimal in Zeno's sequences, so do the associated time intervals. As a result, the description is simply that motion occurs through a series of infinitesimal displacements that occur over infinitesimal time intervals. Had calculus existed in 500 BCE, Zeno would have been laughed out of the room; but it hadn't, and the concept of eternalism that Zeno's paradoxes helped support was able to gain traction.

Regarding Wells' Time Traveller, the initial fallacy is more specific, and again happened right at the outset of his explanation. He began \cite{Wells1895}, ``I do not mean to ask you to accept anything without reasonable ground for it. You will soon admit as much as I need from you. You know of course that a mathematical line, a line of thickness \textit{nil}, has no real existence. They taught you that? Neither has a mathematical plane. These things are mere abstractions\ldots. Nor, having only length, breadth, and thickness, can a cube have a real existence.''

``There I object,'' said Filby. ``Of course a solid body may exist. All real things----''

``So most people think. But wait a moment. Can an instantaneous cube exist?\ldots Can a cube that does not last for any time at all have a real existence?\ldots Clearly\ldots any real body must have extension in four dimensions: it must have Length, Breadth, Thickness, and---Duration. But through a natural infirmity of the flesh, which I will explain to you in a moment, we incline to overlook this fact. There are really four dimensions, three which we call the three planes of Space, and a fourth, Time. There is, however, a tendency to draw an unreal distinction between the former three dimensions and the latter, because it happens that our consciousness moves intermittently in one direction along the latter from the beginning to the end of our lives.''

With respect to our above definitions, the initial fallacy is introduced when the Time Traveller claims that a cube can have no real existence, \textit{and then} equates this to the idea that an instantaneous cube ``exists,'' which is clearly false, so that he suggests alternatively that real bodies are four dimensional. Here, Wells has conflated the concept of an existing object with an instantaneous one. An \textit{instant} does not \textit{exist}, it \textit{happens}, so an instantaneous cube does not exist any more than the instant you reach the period at the end of this sentence exists. A cube can exist, but an instantaneous cube iz three-dimensional, and the moment you reached that last period iz a moment in time. And neither can a four-dimensional world-tube, with length, width, breadth and duration, \textit{exist}---at least, not without conferring an \textit{additional} meta-time dimension in which the world-tube is thought to exist. 

A cube that exists is a cube that endures for some length of time. That existence is a physical property of the cube, but it does not mean that the cube must be a world-tube that extends into the past or future; only that this three-dimensional thing exists in the existing three-dimensional world, for some duration. Otherwise, if the cube iz truly instantaneous, then it just \textit{happens}, and you could say ``a cube momentarily happened over there and then vanished just as quickly as it showed up.'' 

It was by conferring existence on four-dimensional objects that Wells was then able to weave a time travel story, in which a person hops around through a meta-existing space-time, remaining a while and existing there, etc---much like The Doctor exists, from the start to the end of a \textit{Doctor Who} episode, travelling to different regions of a meta-existing space-time, existing and having experiences at each space-time location, altering events there-then in subtle ways, so that the meta-existing space-time subtly changes from the start to the end of an episode. In contrast, a truly four-dimensional space-time with no added sense of existence would not make for much of a tale.

Finally, regarding \citeauthor{Greene2004}'s~\citeyearpar{Greene2004} discussion of the consequences of the relativity of simultaneity, the problem with these descriptions is again the inappropriate use of verbs to describe things that do not \textit{exist} in a four-dimensional theory of relativity. For example, ``The only thing that's real is the whole of spacetime,'' in which the use of ``is'' does confer a sense of existence on spacetime, through the suggestion that it \textit{is} real. The following paragraph then makes this concept much more explicit, as he notes that all the events of space time ``\textit{are}. They all exist. They eternally occupy their particular point in spacetime'', etc. For here, the eternity he is speaking of is not the one that extends from the infinite past to the infinite future of space-time. Here, he explicitly invokes an additional meta-time dimension, stating that all of space-time \textit{meta-exists}, frozen for all \textit{meta-eternity}. 

Proponents of the block universe theory \textit{do} commonly conceive of space-time as existing, in the same sense as the three-dimensional universe \textit{exists}---\textit{now}---according to the above definitions. And just as an existing three-dimensional universe requires four dimensions for a complete description of events that occur in it as it exists, a complete physical description of an \textit{existing} block universe requires a fifth dimension to capture that sense of meta-existence---even if the block \textit{is} forever frozen and unchanging, though more clearly so in related proposals that the block universe evolves and grows with the present \cite{Ellis2014}. The alternative is to embrace and fully explore, from physical, philosophical, and psychological perspectives, the concept of a four-dimensional block universe that merely \textit{iz}---that explicitly \textit{does not} exist---and whether that can be reconciled with all empirical evidence.

Technically speaking, our language \textit{is} already sufficient to explore the concept of a purely four-dimensional block universe. In English, it's possible to structure sentences so they avoid making space-time the subject of a sentence and applying a copula to it. However, unless one is extremely careful, in order to avoid literal inaccuracies that occur in the \textit{regular} use of language the non-existential copulas ``iz'' and ``zare'' were invented above, which explicitly negate the implication of a subject's existence. There are other languages that either lack a copula entirely, or use it more flexibly so that the verb-like quality of ``is'' in English is reduced; but in these cases the sense of the subject's \textit{existence} still tends to be implied. In these cases, instead of avoiding the issue of an explicitly existential copula, the problem actually becomes more subtle.

In any case, the problem of describing existence using languages that either explicitly or implicitly confer existence on the subject of a sentence through basic sentence structure is evident. Furthermore, it is challenging to even think about a reality that strictly does not exist, and this difficulty is reflected in the way that language has developed. As a result, the literal inaccuracies in descriptions like \citeauthor{Einstein1938}'s~\citeyearpar{Einstein1938} and \citeauthor{Greene2004}'s~\citeyearpar{Greene2004} seem to stem from a combination of underdeveloped language and the fact that the theory they were expounding is actually unrealistic. Not only has the situation been similar to the problem people once had, of trying to describe the colour of the sky or the sea without a word for ``blue'', but compounded with that is the problem that common language tends to mean the opposite---our languages tend to be focused around describing things as being or existing---and the reason for this is that existence does seem to be such a fundamental aspect of reality that it is very difficult to force our minds to make sense of a reality that explicitly does not exist. In that sense, the situation is similar to the word ``particle,'' which every particle physicist knows is a very misleading term to use according to modern physics \cite{strassler2024}. Yet the difference here is that whereas particle physicists will acknowledge that ``particle'' is a misnomer, many people who discuss the nature of space-time according to relativity do not appreciate the second sense of existence they often embed quite explicitly---\textit{emphatically}, even---in their descriptions. 

And this problem---the one of space-time's supposed existence; the literal and conceptual error that it is---needs to be eliminated so that we may rectify misunderstandings that have become common in discussions of the meaning of relativity. It is important that proponents of the block universe theory stop misusing verbs like ``is'' and ``exist'' when describing their favoured theory; that they stop misunderstanding the four-dimensional physical world they consider most objectively supported by the cold hard facts of modern physics; that they stop consoling themselves by seeking refuge in the perceptions of the human mind---as any sense of existence, frozen or not, belies an additional sense of time that has never been formalised as part of the mathematical theory. Relativity does not describe space-time as something that \textit{exists}. 

But without this sense of existence, a relativistic block universe is very difficult to reconcile with our experience of the world we live in. As \citet{Greene2004} put it, ``if you agree that your now is no more valid than the now of someone located far away in space who can move freely, then reality encompasses all of the events of spacetime.'' Here, he has not made the mistake of saying that spacetime ``exists'', and his description is consistent with a notion that all the events of spacetime must happen at-once. That is the only consistent explanation unless we add another temporal dimension to the concept and to our physical theory, and describe an existing---enduring---space-time. A block universe that simply iz, which does not exist or endure through an added sense of time is impossible to reconcile with our experience with the physical world. Indeed, as \citet{Greene2004} noted, Rudolph Carnap recounted a conversation about this very topic in which ``Einstein said that the problem of the now worried him seriously. He explained that the experience of the now means something special for man, something essentially different from the past and the future, but that this important difference does not and cannot occur within physics. That this experience cannot be grasped by science seemed to him a matter of painful but inevitable resignation.'' Or, as Greene himself put it, ``it's hard to maintain full resistance to the\ldots viewpoint [that science is unable to grasp a fundamental quality of time that the human mind embraces as readily as the lungs take in the air]. Time is a subtle subject and we are far from understanding it fully. It is possible that some insightful person will one day devise a new way of looking at time and reveal a bona fide physical foundation for a time that flows. Then again, the discussion above, based on logic and relativity, may turn out to be the full story. Certainly, though, the feeling that time flows is deeply ingrained in our experience and thoroughly pervades our thinking and language. So much so, that we have lapsed, and will continue to lapse, into habitual, colloquial descriptions that refer to a flowing time.'' 

This misuse underscores the value of distinguishing between existential and non-existential copulas like ``iz'' and ``zare,'' which clarify that events in space-time, as well as the entire block universe, merely ``happen'' rather than ``exist'' in the enduring sense.

I have attempted here to clarify the distinction between a purely four-dimensional block universe, one that does not exist, does not endure, and should not be conceived to do so by physicists arguing that it iz the most objective outcome based on logic and relativity. That language obscures the formal implication and confounds an accurate conception of the result. The language developed in this section is intended to help us move this discussion forward, using words---most critically copulas---that explicitly do or do not confer a sense of existence---of endurance---on the subjects we are describing.

As I've said, I believe \cite{Greene2004}---and \cite{Einstein1938}---are wrong in their conclusion that a true, four-dimensional, non-enduring block universe iz the most objective interpretation to make based on logic and relativity. As I'll argue in the next section, the problem with that deductive argument comes down to an incorrect assumption about simultaneity---Einstein's operational definition of \textit{simultaneity}---which is outdated and inconsistent with the empirical evidence. This sets the stage for re-examining the operational definition of simultaneity and exploring how a cosmological reinterpretation of relativity aligns more closely with both empirical evidence and our clarified concept of existence. By revisiting simultaneity through this clarified lens, we pave the way for a more empirically consistent and conceptually coherent framework for understanding relativity and cosmology.

\section{What about relativity?}\label{sec:relativity}

\subsection{Relativity and simultaneity in local frames}

\citeauthor{Einstein1905a}'s~\citeyearpar{Einstein1905a} groundbreaking paper on relativity opens, in the first section, with the following definition of simultaneity \cite{Einstein1952}:
\begin{quote}
Thus with the help of certain imaginary physical experiments we have settled what is to be understood by synchronous stationary clocks located at different places, and have evidently obtained a definition of `simultaneous,' or `synchronous,' and of `time.' The `time' of an event is that which is given simultaneously with the event by a stationary clock located at the place of the event, this clock being synchronous, and indeed synchronous for all time determinations, with a specified stationary clock.
\end{quote}
In the second section, he then discussed the relativity of lengths and times, eventually concluding with a note that
\begin{quote}
we cannot attach any absolute signification to the concept of simultaneity, but that two events which, viewed from a system of co-ordinates, are simultaneous, can no longer be looked upon as simultaneous events when envisaged from a system which is in motion relatively to that system.
\end{quote}

This operational definition of simultaneity marked a significant departure from earlier notions of absolute time. By tying simultaneity to the synchronization of clocks within a reference frame, Einstein established a framework that would underpin his revolutionary ideas about space and time. However, this definition also carried implications about the nature of reality that continue to provoke debate.

Einstein's definition is particularly compelling when applied to local phenomena, such as the relative motion of two electrons. The resulting consequence---that there is no objective, absolute definition of simultaneity---is intriguing, even tantalising. In the train example from \S~\ref{sec:block}, we might ask who has the correct notion of simultaneity: the observers on the train or the stationary onlooker? It seems arbitrary to privilege one frame over the other; e.g., to claim that the Earth's reference frame is the ``true'' rest frame, so that the light ``truly'' reaches the back of the car before the front, even though the two events appear simultaneous in the train's frame.

That said, taking the Earth's rest frame as ``true'' would not necessarily contradict the observations. Anyone who has ever ridden a train has likely noticed that it is possible to mentally flip perspectives, imagining either the Earth rolling past a stationary train or the train moving along a fixed Earth. This often happens unconsciously, for instance, when a train coasts before its final stop, giving the impression that the platform is moving before one realises the train is still in motion.

This thought experiment underscores that observers need not insist on being ``at rest'' within their chosen reference frame. If another frame provides a more empirically natural rest perspective, it can be adopted instead. For example, observers on a moving train can describe themselves as moving and reconcile their observations with the fact that light from a bulb reaches the back of the car before it reaches the front in an external rest frame.  In this reference frame, the events that are described as being ``simultaneous'' are not synchronous; in this case, the light from the bulb reaches the front and back of the train car at the same time---i.e. synchronously---but the observers who share a notion of simultaneity with the outside world conclude that in their rest frame \textit{simultaneous} events are not \textit{synchronous}. This subtle but critical distinction highlights that Einstein's operational definition of simultaneity, equating it with synchronicity, is not universally necessary.

Elsewhere, I have discussed an alternative interpretation where \textit{simultaneous} events are not \textit{a priori} defined as those occurring \textit{synchronously} in any particular frame of reference \cite{Janzen2015}. In the present discussion, I maintain that the ``most objective'' definition of simultaneity should not force it to be synonymous with synchronicity. This principle will guide our exploration of simultaneity as it applies beyond local frames, particularly in the cosmological context.

\subsection{Cosmology and the emergence of a cosmic frame}

At the dawn of the 20th century, when Einstein developed the foundational ideas of special relativity, the universe was thought to consist of stars scattered through an unchanging, static void. The notion that the spiral nebulae were distant galaxies, the concept of cosmic expansion, and the framework of general relativity had yet to emerge.

Einstein's definition of simultaneity made sense in this simpler, static view of the cosmos. However, as our understanding of the universe expanded to include its dynamic and large-scale structure, the assumptions underpinning this definition warrant re-examination. The world in which Einstein formulated relativity offered no compelling reason to privilege one star's motion over another. His operational definition of simultaneity was an elegant solution for a universe thought to lack an overarching structure. It grounded his revolutionary theory of space and time, producing intriguing consequences like time dilation, length contraction, and $E=mc^2$, which have all been empirically confirmed.

Over the years, our cosmological perspective has changed significantly. Let us consider a thought experiment related to Einstein's ``happiest thought,'' which occurred to him in 1907, but we'll add a modern twist. In 1907, Einstein realised the principle of equivalence, which led to the formulation of general relativity theory. His happiest thought was essentially that if a person were to jump from a tall building and throw a rock laterally away from themself, since both the falling observer and the rock would be accelerated at the same rate through Earth's gravitational field, from the falling observer's perspective the rock's path should be described as straight line motion with constant speed, and the situation could be described by special relativity. In fact, the rock doesn't even need to be thrown laterally since all the acceleration due to gravity is roughly constant and can therefore be ignored from the observer's relative perspective. The greater inference is that free fall motion within a gravitational field is inertial, and local descriptions can therefore be made according to special relativity.

Let us take this thought experiment a step further, and consider a situation in which you and a friend are skydiving. You jump together from a plane, holding hands, and at some point you push off from one another. Neglecting air resistance, from your perspective, you can describe yourself at rest while your friend moves away at constant speed. Your friend can describe themself as being at rest while you move away at constant speed. Or you might describe the situation from the centre of mass perspective, in which you both move away from the location of the push at constant speeds. 

The fact that any of these descriptions is just as valid as the next is truly powerful. Relativity means that we can do physics in local reference frames without needing to reference any motion with respect to surroundings, nor any gravitational field while in free fall.

But that is only the start of the beauty of it. While relativity's elegance lies in its ability to describe local frames independently of external motion, this independence does not imply that no overarching frame exists. On cosmological scales, evidence for such a frame is both compelling and unavoidable. When you and your friend jump from the plane, in addition to accelerating towards the ground you also continue on with whatever inertia you had from the plane's motion relative to the ground, and because you only need to worry about relative motion you can just ignore this lateral motion as well. You can also ignore the fact that the Earth is rapidly spinning, at a rate of up to half a kilometre per second if you are near the equator. And you can ignore the facts that the Earth is orbiting the Sun at roughly 30~km/s, that the Sun is orbiting our Galaxy at 230~km/s, and that our Galaxy is moving at 627~km/s through the CMB rest frame.

The truly amazing thing about relativity is not simply that it allows us to ignore all of this group motion when describing the relative motion between terrestrial objects, full stop; it's the fact that we get to do physics in local reference frames in spite of all this \textit{real motion} we \textit{do} have within our universe. It's incredibly powerful, and incredibly simplifying and convenient, but even so this fact does not justify assuming that there is no cosmological state of rest. And the fact that all the development in cosmology this past century indicates, very strongly, that there \textit{is} a cosmic state of rest gives cause to reconsider the basic meaning of relativity and the most objective definition of simultaneity. 

As early as 1920, Eddington recognised the potential emergence of a cosmic frame in a theory that extends beyond local phenomena. Writing in this context, he noted \cite{Eddington1920}: 
\begin{quote}
The world taken as a whole has one direction in which it is not curved; that direction gives a kind of absolute time distinct from space. Relativity is reduced to a local phenomenon; and although this is quite sufficient for the theory hitherto described, we are inclined to look on the limitation rather grudgingly. But we have already urged that the relativity theory is not concerned to deny the possibility of an absolute time, but to deny that it is concerned in any experimental knowledge yet found; and it need not perturb us if the conception of absolute time turns up in a new form in a theory of phenomena on a cosmical scale, as to which no experimental knowledge is yet available. Just as each limited observer has his own particular separation of space and time, so a being coexistive with the world might well have a special separation of space and time natural to him. It is the time for this being that is here dignified by the title `absolute'.
\end{quote}

Eddington's ``absolute time'' foreshadows the modern discovery of a cosmic rest frame. His language, describing this concept as a limitation to be viewed ``grudgingly,'' so that one might be ``perturbed'' by it, highlights the personal biases that physicists bring to such questions. Nearly a century later, this anticipation of a universal frame has been borne out by precise measurements of the CMB, which provides direct evidence of a preferred state of rest. To appreciate its significance, we must explore the evidence in detail.

\subsection{The scientific case for absolute time}

When you look at the world around you, you perceive a three-dimensional reality in which you exist. However, what you actually observe is the light that scattered off objects or was emitted by them some finite time in the past. This light travels through space at a finite speed before reaching your eyes. Light moves quickly, so we see our immediate surroundings almost as they are \textit{now}, but images of more distant objects show them as they were further into the past. For example, we see the Sun as it was eight minutes ago, and we see the near side of the Andromeda galaxy as it was roughly 2.5 million years ago, while the light we now see from its far side was emitted roughly one hundred thousand years before that. This everyday experience of perceiving the past through light underscores a fundamental challenge in observational cosmology: reconciling what we see with what actually exists now. Scientific inquiry into this question has revealed profound insights into the nature of time and the structure of the universe.

The relativity of simultaneity complicates this, especially under Einstein's operational definition, which ties simultaneous events to their synchronicity within a given reference frame. This definition, and the associated inference that whether events can be said to occur ``simultaneously'' depends on one's frame of reference, is motivated by two principles in particular: scientific scepticism, and the principle of locality. The former acknowledges that absolute Truth is forever obscured from us, forcing us to make assumptions about the nature of Physical Reality and deduce consequences about observable phenomena we can compare to and potentially falsify our theories. One such assumption is that we live in an enduring three-dimensional universe. Due to the finite speed of light, we cannot definitively confirm that such a universe exists ``out there''; we can only observe events from the past. If this assumption is true, relativity implies that events occurring simultaneously will not be synchronous in all reference frames. Furthermore, if synchronous events in all reference frames are truly simultaneous, this leads to the conclusion that Reality iz a Block Universe.

The principle of locality motivates the operational definition of simultaneity in a different way. This principle states that an object should only be influenced by events within its past light cone. From this perspective, the existence of a well-defined three-dimensional universe might seem irrelevant to present observations since it lies beyond direct scientific measurement. The operational definition of simultaneity thus appears natural, as it prioritises local measurement. In contrast, the concept of an enduring three-dimensional universe introduces a kind of non-locality, assuming a spatial structure that exists independently of immediate causal connections.

It is in this context that we must interpret the truly remarkable observation of the CMB. Despite principles such as locality and the operational definition of simultaneity, the CMB is both conventionally and practically understood within the framework of a cosmic reference frame, representing an enduring three-dimensional universe. According to the standard interpretation, the universe emitted a blackbody spectrum uniformly at the cosmic time when it cooled enough for photons to decouple from matter. At that moment, the universe was not perfectly uniform; matter had clumped together to form the seeds of later galaxy formation. Light from denser regions was slightly redshifted, while light from underdense regions was relatively blueshifted. Apart from these minor shifts due to local matter distributions, the released blackbody spectrum was remarkably uniform as the universe transitioned from a plasma.

In contrast to images of nearby objects (e.g., the Sun or Andromeda), which similarly carry information like red and blueshifts that help infer their properties, the CMB photons encode a wealth of additional, highly significant information about the universe. The CMB image we observe carries information about the \textit{integrated} expansion rate of space over the time it took these photons to travel to us. Over the past 13.8 billion years, as these photons moved through expanding space, they were continuously redshifted due to the universe's local expansion rate. Therefore, what is truly remarkable is the isotropy of this cosmological redshift. Across every direction in space, the photons have been redshifted by precisely the same amount. This uniformity reflects a cosmic time that progressed at the same rate everywhere, throughout which space neither expanded faster nor slower in some regions, which would have led to differential blueshifts or redshifts. Instead, the integrated redshift from the decoupling epoch to the present moment occurred uniformly in all directions. This consistency lies at the heart of the standard interpretation.

In addition to this isotropic redshift, the CMB contains crucial anisotropies that provide further insight into the universe's structure and history. The first type of anisotropy arises from temperature fluctuations linked to density variations in the early universe, imprinted at the moment of decoupling. These fluctuations, detected as variations of one part in 25,000 in the blackbody spectrum, remained intact over billions of years, indicating the extreme uniformity of cosmic expansion. Layered on top of this anisotropy, however, is the dipole: a redshift in one direction and a blueshift in the opposite, reflecting our motion relative to the cosmic rest frame. This motion has been precisely measured at $369.82\pm0.11$ km/s \cite{Aghanim2020a}, providing direct evidence of a cosmic state of rest.

The empirical evidence that we live in an enduring three-dimensional universe with a cosmic state of rest does not simply come from the present dipole anisotropy of the CMB. Together with our present state of motion, the fact that the anisotropy signature from inhomogeneities in the matter distribution of the early universe has remained imprinted within the CMB as it travelled through space that apparently expanded uniformly in every direction for the past 13.8 billion years, signifies that this well-defined cosmic state of rest \textit{has existed} throughout that entire period, throughout the observable universe. Thus, the evidence does not simply show that our universe \textit{presently} has a well-defined cosmic state of rest, but that it has always had one---that \textit{space, now} has \textit{always} been well-defined throughout the history of the observable universe. 

Now, while the CMB provides the clearest and most direct evidence that we live in a three-dimensional universe that exists with objective cosmic time, this is by no means the only such evidence. The scientific evidence supporting the inference that this is very likely a basic fact of reality comes from every cosmological measurement we make. In fact, the principle that there is a cosmic time is just as well supported through scientific evidence as the \textit{less basic} inference that our universe is expanding.

The failure of many physicists and philosophers to acknowledge this empirical evidence underscores the need to revisit how science works. Science proceeds by forming hypotheses about the fundamental nature of reality---often concerning phenomena that cannot be directly observed---deducing observable consequences from these hypotheses, and then testing those predictions against empirical data. If observations confirm the predictions, the hypothesis remains viable, though never definitively proven; if data contradict the predictions, the hypothesis is falsified, requiring revision or replacement. Science, therefore, never proves hypotheses true but continually tests them, verifying their consistency with observations over time. If a hypothesis is true, the best that science can do is verify, over and over again, through multiple observations that confirm predictions, that these basic principles of our theories and models appear to be right based on the best evidence at hand.

Now consider the evidence that our universe is expanding. The first clear indication came when \citet{Hubble1929} combined his distance measurements with the redshifts of galaxies observed by \citet{Slipher1913} and others (see, e.g., \citet{Eddington1923} for an early compilation). Hubble found a roughly linear velocity-distance relation, now known as Hubble's law. Initially, there was some confusion about how to interpret this result since people like \citet{Einstein1922} and Eddington had either mindlessly dismissed or ignored the work by people like \citet{Friedmann1922,Friedmann1924}, \citet{Lemaitre1927} and \citet{Robertson1928}, who had been working to explain the redshift measurements via relativistic cosmological models describing an expanding universe \cite{Janzen2014}. However, following Hubble's result, these previous works surfaced and within the next few years a consensus formed that Hubble's law indicates that we live in an expanding universe; e.g.\ \cite{Eddington1933}.

The standard interpretation elegantly incorporates the cosmological principle that Earth does not occupy a privileged position in the universe. Crucially, this model assumes basic principles, including the existence of a cosmic time, isotropy, homogeneity, and synchrony of space in the cosmic rest frame \cite{Robertson1933}. These assumptions underlie the standard cosmological model, which has been rigorously tested and verified through increasingly precise measurements over the past century. Observations of the redshift-distance relation, the CMB, and the large-scale structure of galaxies consistently confirm the model's principles, including its foundational premise that reality is a three-dimensional enduring universe evolving over cosmic time \cite{Aghanim2020a,Riess2023}.

Moreover, additional evidence from galaxy motion within superclusters---gravitationally dominated watersheds in the cosmic landscape---further corroborates the model. Subtracting the universe's expansion reveals the relative motions of galaxies, which align with the predictions of gravitational interactions in a coherent, three-dimensional framework \cite{Tully2014}. Together, these diverse lines of evidence not only support the expansion of the universe but also substantiate the deeper, more basic principle that there is a cosmic time underpinning this evolution.

\subsection{Relativity in light of evidence}

Thus, when relativists (physicists and philosophers alike) perpetuate the claims that a preferred time can never be observed and that the most objective definition of simultaneity is the operational one, and therefore that the most objective interpretation of relativity, logic, and cold hard science is that reality is---though more accurately iz---a block universe, etc., they do so in utter and bewildering denial of the overwhelming scientific evidence to the contrary from the entire field of modern cosmology. It is logically inconsistent to deny that the empirical evidence supports the basic principle that there is a cosmic time and associated cosmic state of rest while at the same time accepting that the evidence supports the assumption that the universe is expanding over the course of this cosmic time. The latter assumption of our scientific model is already predicated on the former, and can't be true on its own. We \textit{have} verified through a variety of measurements, not only the principle that our universe is expanding but also the more basic principle that it does so over the course of a well-defined cosmic time. And with the CMB measurements, not only have we confirmed that this expansion is so uniform throughout space and time as to have preserved details imprinted on the CMB when the universe was only a few hundred thousand years old, but it also gives us a precise measurement of our present velocity with respect to the cosmic rest frame. 

Ironically, the evidence that there is indeed a cosmic time is far greater than the evidence that supported the heliocentric hypothesis over geocentrism during the Scientific Revolution, as the evidence already comes from a wide variety of high precision measurements. In contrast, stellar parallax, which provides the clearest and most direct empirical verification that Earth is orbiting the Sun, was not observed until two hundred years \textit{after} the deaths of Kepler and Galileo. Additionally, in contrast to the debate between geocentrism and heliocentrism during the Scientific Revolution, in the field of cosmology itself there is no serious challenge to the standard model's foundational principle that there is a cosmic time. But for some reason, relativists seem intent on ignoring the fact that the success of modern cosmology constitutes scientific evidence in support of our model's basic principles; intent on believing we haven't confirmed---and that we actually can't confirm---what we clearly and overwhelmingly already have done. 

It's science denialism, and we're doing it to ourselves, ignoring our own evidence all while telling ourselves we are being ``the most objective.'' We've gaslit ourselves. When \textit{relativists} claim that the most objective conclusion of relativity, logic and the cold hard facts of science imply ours iz a block universe \cite{Greene2004,Einstein1938}, and \textit{eternalists} promote the block universe concept on the basis that ``there is no objective way to slice up spacetime into instants, so we must not assign fundamental significance to any particular slice'' \cite{Wharton2015}, and even \textit{presentists} admit that there can be no empirical way to distinguish the cosmic rest frame \cite{Zimmerman2011}, it's as if none of them understands how \textit{science} works. More importantly, it's as if they don't understand \textit{why} science \textit{must} work the way it does. They've forgotten that science is, at its core, an exercise in maintaining a healthy sense of suspended disbelief in our basic assumptions about reality while working to verify those assumptions in a variety of ways, or else falsify them. We've known for millennia that we have no access to the things that exist---not in and of themselves, anyway---but to the signals that they emit; we see the \textit{phenomena}, not the \textit{noumena}. We observe the things that happen, after they happened, not the things that \textit{are}. We know this---or at least every philosopher of science, and every physicist who claims to be a scientist \textit{should} know it---and therefore must maintain a healthy sense of scepticism towards our hypotheses as we develop our theories and work out new ways of verifying or falsifying them. When exploring the foundations of physics, we allow our curiosity about reality to guide us as we seek to understand the world despite knowing we can never be sure we've got it right. We continually refine both our knowledge and our understanding, working from what we know, what we've determined through empirical observation, to infer the most consistent interpretation of the facts we observe---and then working from the picture of the world that we think we understand, to think up new questions that we can go out and investigate.

And when it comes to the theory of relativity, we've found that the \textit{empirical evidence} from cosmology, combined with the basic logical structure of science, strongly indicates that which had been thought at the turn of the twentieth century to be impossible: from a \textit{scientific} perspective, modern cosmology indicates that there is a cosmic state of rest, an objectively well-defined \textit{now} that our observations have not falsified. We've determined this indirectly through the success of our expanding universe models which are predicated on the assumption that there is a cosmic time, as well as through direct measurement of our velocity through the universe, through the CMB dipole. We've found that our universe expanded to its present state from an initial hot big bang 13.8 billion years ago (though this is not necessarily the age of particles in local frames due to time dilation effects), with the terms defined in \S~\ref{sec:existence} applying very well to describe the \textit{existence} of our three-dimensional \textit{universe} and the \textit{occurrence} of \textit{events} that form four-dimensional \textit{space-time}. In contrast, it would be very difficult to reconcile a theory that accounts for all the cosmological observations, particularly the CMB features, without assuming an absolute cosmic time and uniform expansion throughout the universe, throughout the course of that cosmic time \textit{a priori}---and currently no serious challenge to this interpretation exists.

In fact, the main challenge seems to come from relativity itself, interpreted in ignorance of the current cosmological evidence, and taken to support a vague concept that reality \textit{is} four-dimensional, that four-dimensional space-time \textit{exists}, which is literally and conceptually inaccurate---a false interpretation of the physical theory that comes from allowing an unjustified and formally unaccounted for tacit concept of existence to sneak into the general idea undetected. This idea of four-dimensional existence then enables us to explore such tantalising possibilities as time travel \cite{Morris1988,Hawking1992,Thorne1994,Visser1995}, the emergence of space-time from quantum phenomena \cite{Carroll2019}, that reality might be an evolving block universe \cite{Ellis2014}, or simply the intriguing challenge of reconciling our tacit experience that time passes with the notion that time is simply an illusion \cite{barbour2001,davies2002,Greene2004,callender2017,rovelli2019}. Formalisation of that tacit concept of existence, as discussed in \S~\ref{sec:existence}, clearly shows that each of these ideas is fundamentally flawed, as they are all predicated on an underlying concept of existence that is distinct from the time-dimension within space-time. 

Now, regarding the separate issue of locality, we also note that the continuously evolving rate of expansion along any particular line of sight has been exactly the same as the evolving rate observed along every other line of sight, even though those regions have not been causally connected since the CMB was emitted in the early universe. Without a prior mechanism such as the hypothesised inflationary epoch \cite{Guth1981} that could mean diametrically opposite regions of the observable universe were causally connected earlier on, it would appear that this three-dimensional universe we find ourselves existing in is highly nonlocal; that its expansion has emerged through a mechanism that acts the same everywhere, potentially without causal connection to explain the apparent uniformity. 

In spite of the legitimate concerns that initially motivated adopting the operational definition of simultaneity---the notion that we might not be able to observe our motion through space, that adopting a preferred state of rest in our universe adds non-objective, non-egalitarian structure to a relativistic description of reality, and that this structure amounts to a reality that is in essence nonlocal---the \textit{empirical evidence} has shown us that those initial concerns were \textit{actually} unwarranted---that we can indeed empirically determine, with great precision, the preferred state of rest in a universe that does exist beyond our past light cone. 

Thus, just as you and your friend on the train in the thought experiment above are not limited to an operational definition of simultaneity when you have access to observations of your surrounding world, so are we, here on Earth, not limited to Einstein's definition of simultaneity and its associated ontological interpretation that there is no objective state of rest. On the train, you can certainly still synchronise clocks and define simultaneity by the set of events that occur synchronously in your local rest frame. But by referencing what you observe in your surroundings---that you are on a train rolling along the ground---you naturally infer a different state of rest and form an associated definition of simultaneity that is different from the operational one. According to that definition, the events that occur at the same time on your synchronised clocks do not occur ``simultaneously.'' Similarly, by referencing cosmological observations we do infer an objective state of rest for our entire universe. The empirical data indicate that there has been a cosmic state of rest that has existed for 13.8 billion years since before the CMB was released from the primordial plasma, which has existed throughout that time, all throughout the 13.8 billion light-years of expanding space from which we now receive empirical data. 

It's true that we can, here on Earth, synchronise clocks and define a local state of rest, and we can do physics here without reference to the external universe. Relativity allows this. But to define synchronised events in this frame as simultaneous, and to claim that \textit{this} state of rest is \textit{the most} objective one for us to use, amounts to willful denial of significant and highly relevant empirical evidence to the contrary. In light of that evidence, not only is an operational definition of simultaneity ``not the only option available to us,'' but the evidence actually clearly stands against an egalitarian ontological interpretation of relativity. According to all the evidence available to us from cosmology, there \textit{is} an absolute cosmological rest frame and an absolute simultaneity of events that occur in our enduring three-dimensional universe that do not occur synchronously in general frames of reference. 

\section{What about black holes?}\label{sec:BH}

\subsection{Gravitational collapse and the nature of black holes}

As long as the universe is smooth throughout, it must be possible to define a three-dimensional general relativistic space that exists, or endures, and through that existence forms four-dimensional space-time as the set of all events that ever happen in it. For example, consider a large trampoline, with bowling balls of varying densities rolling around on it, growing and shrinking and merging together. The denser these balls become, the more the trampoline warps under their concentrated weight. 

By this conception, the two-dimensional trampoline surface represents a slice of three-dimensional space that might be thought to warp in the presence of gravitational mass, as described by general relativity. And just as the two-dimensional trampoline exists in our world, with individual moments having different well-defined states that all unfold to create the tapestry of events that is space-time, so too can the existence of a three-dimensional space be thought of in this way. Locally, in the presence of large or dense mass, not only should space be perceived to have curvature, but the apparent passage of time may be affected as well, so that space-time can be said to be curved by mass. As long as the curvature of space is smooth, the trampoline always remains diffeomorphic to a simple flat plane, and if the apparent local passage of time at different points of the trampoline maps smoothly to the global passage of time in which the trampoline ``exists,'' then clearly the cosmological slicing of space-time can be objectively defined.

But what about singularities? What if one of the bowling balls shrinks all the way down to a single point, and the trampoline becomes infinitely curved there? This would indeed be a problem, as it would no longer be possible to describe that as a point in ``space,'' and its ``existence'' would be impossible to define. Therefore, we must consider whether it is possible to justify the development of space-time singularities through gravitational collapse in a cosmologically relativistic universe---i.e. one in which an absolute, cosmic simultaneity is defined \textit{a priori}, for the reasons discussed above. 

Let us consider the local space-time metric outside a non-rotating, uncharged, spherically-symmetric body that exists in asymptotically flat, three-dimensional space. Defining the radial coordinate as $r$ and the time coordinate as $t$, and starting with the most general functions of $r$ and $t$ for the metric components of a space-time geometry with radial symmetry about $r=0$, we note the significant result first derived independently by \citet{Jebsen1921} and \citet{Birkhoff1923}. The Jebsen-Birkhoff theorem states that, without loss of generality, the geometry of general-relativistic space-time with this specified general form of the line-element is described by the stationary Schwarzschild line-element,
\begin{equation}ds^2=-\left(1-\frac{r_{\mathrm{Sch}}}{r}\right)dt^2+\frac{1}{1-\frac{r_{\mathrm{Sch}}}{r}}dr^2+r^2(d\theta^2+\sin^2\theta d\phi^2),\label{eq:Sch}\end{equation}
where $d\theta^2+\sin^2\theta d\phi^2$ describes the two spherically-symmetric dimensions of space, and $r_{\mathrm{Sch}}=2GM/c^2$ is the radius of a coordinate singularity---the ``event horizon''---which depends on the mass $M$ of the central body, the gravitational constant $G$, and the speed of light $c$. Remarkably, the metric components turn out to be independent of the specified time coordinate $t$.

The causal structure of this geometry in two different coordinate frames is shown in Figure~\ref{fig:BH}, where the spherically-symmetric part of space has been suppressed, and ``infalling'' paths of photons moving purely through $t$ and $r$ are plotted as blue curves, while ``outgoing'' photons moving purely through $t$ and $r$ are plotted as red curves. Figure~\ref{fig:BH} shows paths of these radially moving photons, both in the stationary Schwarzschild coordinates of Equation~\ref{eq:Sch} (Figure~\ref{fig:BH}a) and in Eddington-Finkelstein coordinates (Figure~\ref{fig:BH}b). The Eddington-Finkelstein coordinates are found by introducing a variable $v=t+r+r_{\mathrm{Sch}}\ln|r/r_{\mathrm{Sch}}-1|$, which runs in the direction of the infalling (blue) photon geodesics, and setting $v-r$ as the vertical ``time'' axis so that, when plotted against $r$, infalling photon paths are straight 45-degree lines. The main purpose of the Eddington-Finkelstein diagram is to clearly show that both infalling photon lines and infalling timelike lines extend beyond the event horizon and reach the singularity that forms at $r=0$ in finite proper time. Thus, as opposed to Figure~\ref{fig:BH}a, in which timelike worldlines (such as that which might represent the surface of a collapsing star) appear to approach the event horizon asymptotically as the coordinate $t\rightarrow\infty$, Figure~\ref{fig:BH}b clearly shows that infalling worldlines will pass smoothly through the event horizon and all the way to $r=0$ in finite proper time.

In Figure~\ref{fig:BH}a, both the ``infalling'' and ``outgoing'' photon paths asymptotically become vertical at the event horizon, $r=r_{\mathrm{Sch}}$, whereas only the ``outgoing'' photon paths become vertical at $r=r_{\mathrm{Sch}}$ in the Eddingron-Finkelstein diagram, Figure~\ref{fig:BH}b. I'm calling these ``infalling'' and ``outgoing'' photon paths, with quotation marks, because of the change in causal structure of the geometry at $r=r_{\mathrm{Sch}}$: as is clear from Equation~\ref{eq:Sch}, when $r<r_{\mathrm{Sch}}$ the coefficient of $dt^2$ has an overall positive sign and the coefficient of $dr^2$ is negative, since $r_{\mathrm{Sch}}/r>1$ in this interval; therefore, when $r<r_{\mathrm{Sch}}$, $r$ is timelike and $t$ is spacelike. Therefore, as shown in both Figure~\ref{fig:BH}a and Figure~\ref{fig:BH}b, at $r<r_{\mathrm{Sch}}$ both sets of photons follow paths moving towards $r=0$, including the ``outgoing'' ones. In that region, in the stationary Schwarzschild frame shown in Figure~\ref{fig:BH}a, a massive particle that must always move slower than the speed of light can move either up or down in $t$, but cannot cross any blue light paths below or any red light paths above. Within the Eddington-Finkelstein diagram shown in Figure~\ref{fig:BH}b, a particle moving slower than the speed of light cannot pass below or to the left of the blue light paths and cannot pass above or to the right of the red light paths. Thus, in both cases at $r<r_{\mathrm{Sch}}$ particles that move slower than the speed of light must evolve towards $r=0$. On the other hand, at $r>r_{\mathrm{Sch}}$, particles moving more slowly than the speed of light can make their way either towards $r=r_{\mathrm{Sch}}$ or outward in the direction of ever increasing $r$.

\begin{figure*}[t]
\centering
\includegraphics[width=3.3in]{Sch.png}\includegraphics[width=3.3in]{EF.png}\\
\hspace{0.1in}(a)\hspace{3.05in}(b)
\caption{Causal structure of the Schwarzschild geometry, shown in (a) stationary Schwarzschild coordinates and (b) Eddington-Finkelstein coordinates. The spherically symmetric part of the geometry is suppressed in both diagrams, which therefore show worldlines of photons moving purely radially ``inward'' (blue) and ``outward'' (red). Lightlike lines bunch to form a solid vertical line at the event horizon, $r=r_{\mathrm{Sch}}$. At $r<r_{\mathrm{Sch}}$, $r$ is a timelike coordinate, and both light and matter in this region must evolve towards $r=0$.}
\label{fig:BH}
\end{figure*}

Given that the surface of a star must move more slowly than the speed of light, let us now consider the physical description of spherically-symmetric gravitational collapse of a mass with initial radius $r>r_{\mathrm{Sch}}$, outside of which the space-time geometry is given by the Schwarzschild line-element, Equation~\ref{eq:Sch}. 

As described in the stationary Schwarzschild coordinate system and shown in Figure~\ref{fig:BH}a, the surface of the star can only reach the event horizon as $t\rightarrow\infty$, since the infalling (blue) light lines become vertical as $r\rightarrow r_{\mathrm{Sch}}$. On the other hand, as shown in Figure~\ref{fig:BH}b, the surface of the star can follow a finite timelike path, never moving left through any blue line or right through any red line, so that it must reach $r=0$ a finite interval after crossing the horizon. In fact, as proven by \citet{Penrose1965}, the star can actually be rotating and need not be perfectly spherically symmetric, as assumed here, in order for it to form a singularity in finite proper time: any closed trapped surface, such as the misshapen surface of a star that has crossed the event horizon, must collapse to a singularity within a finite proper time interval.

So, if this is all true, how can we reconcile the problem noted above, that in order for there to \textit{be} a well-defined, three-dimensional \textit{existing} universe, it cannot develop singularities? The key is to be careful and precise about \textit{when} the star's surface can actually be said to reach the event horizon, from the perspective of external observers (like us). And to do that, \textit{we must ensure an accurate definition of} cosmic simultaneity \textit{that is independent of the coordinate system being used}. 

Let us start by considering the collapse scenario from the perspective of a local observer who describes the star as an existing spherically-symmetric mass. From that local perspective, space-time is described by Equation~\ref{eq:Sch}, in which $r$ is by definition the radial coordinate of spherically symmetric space-time, and $t$ is the associated time variable in which spherically symmetric space was defined to exist. Therefore, as the star's surface collapses towards $r=r_{\mathrm{Sch}}$, reaching the event horizon in finite proper time on its own clock, from the perspective of external observers in the local reference frame, watching a star collapse, its radius should always be given by whatever value of $r$ it has reached along lines of constant $t$---i.e. at times given by horizontal lines in Figure~\ref{fig:BH}a, according to which the star asymptotically approaches $r=r_{\mathrm{Sch}}$ as $t\rightarrow\infty$.

From this perspective, the star never reaches its event horizon, never evolves along timelike $r<r_{\mathrm{Sch}}$, and never forms a singularity. The black hole described from the perspective of this observer, who remains at constant radius outside the collapsing star and describes the external space-time geometry as an enduring, spherically-symmetric space, is a very different object from the standard conceptual model of a black hole. From the perspective of the external observer who uses stationary Schwarzschild coordinates, the star \textit{remains forever} outside its event horizon, its contents causally connected to the rest of the universe. Matter may continue to fall into this ``black hole,'' adding to its mass, while information about the particles that make up this asymptotically collapsing body remain forever accessible, at least in principle, to the outside world.

In contrast, the standard conceptual model of gravitational collapse to a black hole is more readily inferred from pictures like Figure~\ref{fig:BH}b, in which dynamical collapse occurs over the course of $v-r$, with individual moments described by horizontal slices of constant $v-r$ (see Figure~\ref{fig:BH_EF}b). According to this description, the external observer may remain forever at $r>r_{\mathrm{Sch}}$ while the star's surface crosses the event horizon and reaches $r=0$ in finite $v-r$. Afterwards, the event horizon, timelike $r<r_{\mathrm{Sch}}$, and the singularity at $r=0$ are all described to ``exist'' with spacelike separation from the external observer. The event horizon in this case is a one-directional lightlike surface that matter can fall into, and out of which no information can ever escape. 

\begin{figure*}[t]
\centering
\includegraphics[width=3.3in]{EF_Sch_t.png}\includegraphics[width=3.3in]{EF_v-r.png}\\
\hspace{0.1in}(a)\hspace{3.05in}(b)
\caption{The Schwarzschild geometry displayed in two Eddington-Finkelstein diagrams similar to Figure~\ref{fig:BH}b, shown here with black lines depicting (a) slices of constant $t$ (defined at at $r>r_{Sch}$) and (b) slices of constant $v-r$.}
\label{fig:BH_EF}
\end{figure*}

According to the latter interpretation, one may claim, e.g. \cite{Penrose1969,penrose2002}: ``No matter how long the external observer waits, he can always (in principle) still see the surface of the star as it \textit{was} just before it plunged through the Schwarzschild radius.'' Indeed, the emphasis on ``\textit{was}'' is significant in this sentence, as it conveys the meaning that at \textit{later} proper times, from this external observer's perspective, the collapsing star ``has \textit{already} crossed'' its event horizon and ``\textit{formed}'' a singularity at $r=0$.

\textbf{There is a subtle but significant point to be noted when interpreting the gravitational collapse scenario through figures like those shown in Figure~\ref{fig:BH}} (i.e.\ as depicted explicitly in Figure~\ref{fig:BH_EF}b): \textit{general relativity does not specify} a priori \textit{any coincidental proper time evolution along separate worldlines}. In practice, probably following intuition developed according to Einstein's operational definition of simultaneity, we \textit{tend} to interpret synchronous events as coincidental---as simultaneous occurrences---when interpreting space-time diagrams. Thus, e.g., when interpreting Figure~\ref{fig:BH}a it is clear that as the external observer ``remains forever'' at $r>r_{\mathrm{Sch}}$, a collapsing star's surface should asymptotically approach the event horizon. And when interpreting Figure~\ref{fig:BH}b, it is clear that while that external observer ``remains forever'' at $r>r_{\mathrm{Sch}}$, and in principle can forever see photons emitted by the star's surface before plunging through the event horizon (though practically these photons soon become too redshifted to detect), there is a well-defined sense in which the external observer is ``\textit{later}'' seeing the star's surface ``as it \textit{was}, just \textit{before}'' it plunged through the Schwarzschild radius---i.e. according to which the star's surface ``reaches its event horizon coincidentally with some finite proper time on the external observer's clock,'' at synchronous $v-r$, ``after'' which the external observer ``\textit{goes on existing} while the star collapses to form singularity at $r=0$'' in finite $v-r$ and continues to exist ``afterwards.''

Figure~\ref{fig:BH_EF} depicts both of the physical interpretations described in the above paragraph, within the same Eddington-Finkelstein frame of reference, with lines of ``simultaneity'' that define coincidental events across ``space'' drawn in black. \textit{General relativity favours neither picture over the other} a priori, \textit{and both depict valid dynamical sequences of} simultaneous \textit{events when considered purely from the standpoint of the mathematical theory}. In Figure~\ref{fig:BH_EF}a, it is clear that as an external observer remains at some fixed $r>r_{\mathrm{Sch}}$, they can forever watch the collapse of a star until photons become too redshifted to detect, and they can subsequently drop particles, probes, and anything else into the star. \textit{And all of these things will asymptotically approach the event horizon for the rest of time, all finally arriving there} simultaneously. Conversely, Figure~\ref{fig:BH_EF}b corresponds to the standard conceptual model of black hole formation described above.

It is worth remarking that since the 1960s physicists have favoured the latter dynamic gravitational collapse scenario over the former, even though general relativity does not specify which events happen coincidentally along separate worldlines; i.e. which two separate events should be described as happening ``simultaneously'' as proper times advance. In particular, the mathematical theory of relativity provides no prior specification of ``when'', along the external observer's worldline, the star's surface crosses its event horizon. And general relativity says nothing about a section of the external observer's worldline existing ``after'' the star's surface has crossed the event horizon. \textit{This is purely physical interpretation of the mathematical description}. 

One might argue that the foliation in Figure~\ref{fig:BH_EF}a is incomplete, and does not specify the coevolution of matter \textit{after} the event horizon. However, given that the $t$-slicing tilts and becomes vertical at the horizon in this frame, and that $r$ describes the timelike radial direction of spherically symmetric space-time geometry after the horizon, a reasonable guess would be that constant $r$ slices describe coincidentally occurring, simultaneous events, for the rest of time until reaching the final singularity at $r=0$---i.e. that we could reasonably add vertical black lines to Figure~\ref{fig:BH_EF}a at $r<r_{\mathrm{Sch}}$.

While general relativity is agnostic regarding the particular foliation of space-time, or equivalently the coevolution of proper time along separate worldlines, the collapse picture presented in Figure~\ref{fig:BH_EF}a is significantly different, from a physical standpoint, when compared to the standard scenario that has been inferred from figures like Figure~\ref{fig:BH_EF}b. According to the foliation shown in Figure~\ref{fig:BH_EF}a, neither the singularity nor even the event horizon should ever be said to ``exist,'' from an outsider's perspective, after any finite amount of proper time has elapsed. In this case, all matter that ever falls into the collapsing star, everything it ever collides with, etc, all reaches the horizon simultaneously. 

If this scenario represents the true physical nature of black holes in our universe, the implication for physics is significant. For our purposes here, such a scenario would imply that singularities do not \textit{exist}, and we would avoid the complication with defining a global cosmic time variable everywhere and always throughout an enduring three-dimensional universe. But also, since black hole event horizons would not \textit{exist}, Hawking radiation should not be emitted \cite{Hawking1975} and there would be no black hole information paradox \cite{Hawking1976}. And since singularities also would not exist, we would not need to worry about the formation of naked singularities \cite{Penrose1969,Engelhardt2024}. Essentially, six decades of black hole physics will have been based on an incorrect interpretation of the \textit{physics} of gravitational collapse.

\subsection{Implications of coordinate slicing and black hole physics}

The significant difference between these two physical interpretations highlights an important fact: from the standpoint of the mathematical theory of relativity, \citeauthor{Penrose1969}'s~\citeyearpar{Penrose1969} statement about the external observer's continued \textit{existence}, \textit{after} the star had plunged beneath its event horizon, was critically presumptive and unjustified. The interpretation is only accurate with respect to a particular foliation of the space-time, and the physical relevance of that foliation had not been justified. 

To investigate further whether the traditional collapse scenario is physically justified, let us consider the $v-r$-slicing of the Schwarzschild geometry, not just as the synchronous time-variable in the Eddington-Finkelstein frame, but also as a foliation of the geometry within the stationary Schwarzschild frame (Figure~\ref{fig:BH_vr}). The $v-r$-slicing in Figure~\ref{fig:BH_vr}a is particularly illuminating, as it shows the effect of the choice in coordinate transformation: by making the \textit{definition} that the paths of infalling photons should be the $v$-direction in space-time, the asymptote that occurs as infalling photons take infinite $t$ to reach the horizon is effectively ``dragged (infinitely far) down'' at the horizon so that infalling photon trajectories become straight lines in the resulting coordinate system. As described in the local reference frame of spherically-symmetric space-time in which $r$ and $t$ were defined, the new ``time'' direction of the four-dimensional space-time geometry, $v-r$, extends from finite $t$ at $r\rightarrow\infty$ to $t\rightarrow\infty$ at $r=r_{\mathrm{Sch}}$, and then from $t\rightarrow\infty$ at $r=r_{\mathrm{Sch}}$ down again to finite $t$ at $r=0$. Thus, every ``moment'' described in this new reference frame extends from some finite $t$ to infinite $t$ outside the horizon (i.e. across $r>r_{\mathrm{Sch}}$), and then inside, or more accurately \textit{after} the horizon, the same ``moment'' extends across the full timelike interval  $0<r<r_{\mathrm{Sch}}$ and across space from some shrinking finite value of $t$ out to $t\rightarrow\infty$. 

\begin{figure*}[t]
\centering
\includegraphics[width=3.3in]{Sch_v-r.png}\includegraphics[width=3.3in]{EF_v-r.png}\\
\hspace{0.1in}(a)\hspace{3.05in}(b)
\caption{The Schwarzschild geometry displayed with slices of constant $v-r$ shown as black lines in (a) the Schwarzschild frame and (b) the Eddington-Finkelstein frame.}
\label{fig:BH_vr}
\end{figure*}

The asymptote in the Eddington-Finkelstein $v$-direction is a mathematical construct intended to show the smooth transition across the event horizon for worldlines of infalling particles. The resulting coordinate $v-r$ is not a physically meaningful ``direction'' of the space-time, particularly from the perspective of external observers who might want to conclude that fully-formed black holes ``exist now'' in their universe. Rather, it is similar to defining a ``time'' direction in a Robertson-Walker space-time in which individual ``moments'' extend spatially outwards to the big bang in one direction and to the end of time in the other. Such a spacelike hypersurface is a three-dimensional set of events that occur over the course of eternity, so in the course of this new ``time'' variable there is a significant sense in which all of \textit{eternity} is described to ``always exist.'' For example, take the local reference frame of a neutrino that has been travelling through the universe since being created a million years after the big bang: due to time dilation, the amount of elapsed proper time on this neutrino's clock may be only a few years, and if ``space'' is described as the synchronous set of events within the neutrino's frame of reference, then for all time the big bang is continually happening somewhere behind it, while the end of the universe is forever in front of it.

Similarly, the coordinate $v-r$ in the Schwarzschild geometry does not describe physically realistic ``moments'' of time in the reference frame of \textit{any} observer outside the collapsing star. In the case of a star collapsing as it moves through a well-defined three-dimensional enduring universe, an external observer who remains a fixed distance away would not find that the cosmic hypersurfaces are described in their frame of reference by the $v-r$-direction, which forms a 2-sphere at $r=r_{\mathrm{Sch}}$ around the centre of the star where this observer's own coordinate time becomes infinite. Rather, the cosmic hypersurface should be tilted a finite amount in the local reference frame due to the external observer and the star's motion through the universe. And in this case, the effect of the coordinate singularity at $r=r_{\mathrm{Sch}}$ on this finitely tilted ``now'' should effectively be unchanged: just as the local coordinate time $t\rightarrow\infty$ at the Schwarzschild radius, so should cosmic time at $r=r_{\mathrm{Sch}}$ tend to infinity, so that the collapsing star is described to asymptotically approach the event horizon in infinite cosmic time. 

Thus, the physical scenario of dynamic gravitational collapse should be effectively similar to our physical interpretation above of the slicing in Figure~\ref{fig:BH_EF}a, not the interpretation inferred from Figure~\ref{fig:BH_EF}b. Along with providing a resolution to several outstanding problems in black hole physics (Hawking radiation, information paradox, naked singularities), this scenario indicates that black holes should not pose a problem for the cosmological interpretation of relativity discussed in \S~\ref{sec:relativity}. Furthermore, if such asymptotic black holes are what truly exist in our universe, there may be implications for black hole merger calculations that could potentially be important in gravitational wave models.

Finally, this collapse scenario resolves an issue related to the concept of physical existence discussed in \S~\ref{sec:existence}. Due to the definition of the ``time'' coordinate $v-r$ in the Eddington-Finkelstein coordinate frame, as a direction within the Schwarzschild geometry that spans all of $r>0$ as well as infinite t, the collapse scenario inferred from the interpretation that slices of constant $v-r$ represent ``moments'' of time leads to a conception in which not just space, but space-time is imagined to ``exist'' within the event horizon. According to this conception, not only do the three dimensions of space described by $t$, $\theta$, and $\phi$ always ``exist'' throughout the course of $v-r$, but $0<r<r_{\mathrm{Sch}}$ ``exists'' as well. Thus, an external observer can ``exist'' outside a collapsing star, and over the course of ``time'' the star can collapse beneath its event horizon to form a massive singularity at $r=0$. From this ``moment'' on, the singularity, the entire timelike interval $0<r<r_{\mathrm{Sch}}$ (along with three spatial dimensions), and the event horizon, are all thought to ``exist'' within the enduring three-dimensional external universe. ``After'' this black hole has formed, the external observer can toss a brick into it, then wait a little while and toss another one in so that at some later ``moment'' the first brick is at, say, $r=r_{\mathrm{Sch}}/3$ and the second is at $r=2r_{\mathrm{Sch}}/3$. But according to the discussion in \S~\ref{sec:existence}, this notion, that bricks can fall, one after the next, through space-time, as if it's an enduring arena \textit{through} which they can move over the course of ``existence,'' involves a second sense of time---a meta-time through which the four-dimensional black hole is imagined to \textit{exist} as a finite ball of four-dimensional space-time within its event horizon. The concept is essentially similar to one where we might imagine ourselves to exist, in this meta-time sense, at our space-time location here on Earth in 2024, while right \textit{meta-now} there may be a brick \textit{meta-existing} in the year 3024---a year that we'll come to in due course, a thousand years from now, at some future meta-moment when that brick will have similarly evolved, say, to 4024, etc.

As noted above, this interpretation appears to be due to the fact that in the Eddington-Finkelstein frame the time-direction of the Schwarzschild geometry was defined as a direction through space-time that spans all of $r$ as well as infinite $t$. But while such an arbitrary slicing of the four-dimensional geometry is allowable as a mathematical abstraction that, e.g., enables us to describe a smooth transition across the event horizon in finite proper time, it is important to be careful when interpreting the events along these slices as dynamically unfolding within enduring three-dimensional spatial slices. In the latter sense, the physical interpretation of the synchronous direction in the Eddington-Finkelstein coordinate frame as ``space'' outside a collapsing star is unjustified.

This example highlights a critical distinction in our physical interpretation of relativity. Traditionally, whether the separation of events is spacelike, lightlike or timelike has been understood to bear critical significance. For example, if two events zare spacelike separated, then through the operational definition of simultaneity there is thought to be a significant sense in which those two events can be described as simultaneous occurrences in ``space.'' But in order to reconcile relativity as a mathematical description of physical reality that dynamically exists, we are forced to relinquish the ontological interpretation that synchronicity bears any significance on the actual simultaneity of events that occur in the course of this existence. Accordingly, the actual three-dimensional universe is described in any coordinate frame as a particular ``spacelike'' hypersurface, and there is otherwise no significant difference in distinguishing any two events that do \textit{not} lie along that hypersurface as ``spacelike'', ``lightlike'' or ``timelike'' separated. Rather, all such spatiotemporal relations are between events that occur at different cosmic times, and potentially at different points in space. Conversely, under the egalitarian interpretation in which any spacelike hypersurface is said to be as significant as any other, when we do not also admit the logical consequence of a block universe and clearly acknowledge the implication that ``space'' in any particular frame cannot then be thought to dynamically exist, we run into misunderstandings and misinterpretations in which events are thought to dynamically unfold, or regions of space-time are thought to exist. And these interpretations are inconsistent with the logical consequences of the basic assumption we wanted to make.

\section{What implications for cosmology?}\label{sec:cosmology}

\subsection{Assessment of cosmological principles for a cosmologically relativistic universe}\label{sec:cosmo_analysis}

Building on the discussion so far, particularly the treatment of black holes and gravitational collapse in \S~\ref{sec:BH}, it should now be evident that the standard cosmological line-element imposes significant constraints on the treatment of cosmic time, limitations that merit reconsideration. Its limitations suggest that we should not \textit{expect} a general relativistic universe to \textit{be} as we have defined it over the past century.

As noted earlier, Eddington observed that ``the world taken as a whole has one direction in which it is not curved; that direction gives a kind of absolute time distinct from space.'' While his acknowledgment that a cosmological state of rest and an absolute cosmic time may arise naturally within a theory of phenomena on a cosmical scale is remarkable, it has not been widely appreciated in the century since. Perhaps this oversight stems from what Eddington himself described as a ``grudging'' reluctance to accept such a limitation, or from the compartmentalisation of relativistic and cosmological worldviews. Regardless, his insight raises a profound question: why should a general relativistic universe lack curvature in the direction of cosmic time while allowing spatial curvature? The persistence of this presumption likely stems from a constrained understanding of relativity, tied to the assumption that simultaneous events occurring ``now'' from any perspective, even the cosmical one, must always align with synchronous definitions of time.

Consider \citeauthor{Robertson1933}'s~\citeyearpar{Robertson1933} derivation. In his section titled ``Cosmic time. A priori specification of the line element,'' he outlined a general coordinate framework for cosmology based on four conditions: (a) the lines parametrised by the coordinate $t$ should be geodesics along which $t$ measures proper time, chosen such that they represent, as consistently as possible, the mean motion of matter in their vicinity; (b) the space $t = \mathrm{const}$ should be orthogonal to this congruence of geodesics; (c) space should be isotropic; and (d) space should be homogeneous.

Robertson's conditions (a) and (b) are interdependent, with (a) relying on the definition of space in (b), and both depending on motion. Robertson introduced the coordinate $t$ as a ``\textit{cosmic time}'' variable (his emphasis), justified by Weyl's postulate that the universe can be represented as a coherent bundle of geodesics. This postulate, in turn, was motivated by astronomical measurements of the mean motion of matter.

Condition (a) establishes a cosmic time variable defined by comoving geodesics representing the world-lines of observers at rest in the cosmic reference frame---a notion that remains empirically justified. Recognising this assumption as consistent with relativistic principles allows us to clarify the four foundational principles of the standard FLRW model: (a) the coordinate $t$ defines a bundle of comoving geodesics, establishing a cosmological state of rest, cosmic simultaneity, and spatial slices representing ``the universe'' at any moment; (b) the direction $t$ is orthogonal to these slices, ensuring synchronicity of simultaneous events in the cosmic rest frame; (c) the universe is isotropic; and (d) the universe is homogeneous. These principles mirror Robertson's framework.

While condition (a) inherently specifies a cosmic simultaneity, condition (b) introduces an additional assumption---motivated only by an operationally relativistic conception of simultaneity---that simultaneous events occur synchronously in the cosmic rest frame. A general relativistic universe should not \textit{require} synchronicity in the cosmic frame \textit{a priori}. Synchronicity may be possible, but it should not exclude isotropic and homogeneous models that lack this property.

While it should already be clear that the synchronicity assumption is poorly motivated from the perspective of the cosmological interpretation of relativity, since its primary motivation comes from the operational interpretation, this assumption actually highlights a much deeper issue with standard FLRW cosmology. The synchronicity assumption of FLRW cosmology is not merely a mathematical convenience; it is equivalent to assuming that the universe \textit{is}, at any cosmic time, \textit{exactly as it appears to be}---that is, the geometry of the universe throughout cosmic time evolves precisely as described in the synchronous time slicing of space-time in the frame of fundamental observers. While this might initially seem like a natural assumption and indeed serves as a convenient first guess when exploring newly discovered phenomena, the notion that reality is exactly as it appears is profoundly at odds with both general relativity and historical lessons in physics and astronomy.

To see why this assumption is deeply flawed, we can consider examples from both modern physics and historical paradigms that demonstrate the dangers of conflating appearances with reality. The first example we might note is the distinction between the operational and cosmological interpretations of relativity we've been investigating. That simultaneous events are defined to be the synchronous events in any given reference frame is a hallmark characteristic of the former, which I've already argued is at odds with empirical evidence and has led to significant problems with our understanding of concepts like existence and space-time, the two of which are often conflated in the notion that space-time is the fabric of reality. The cosmological interpretation of relativity, on the other hand, hints that reality may not be trivially as it appears, e.g.\ as it tells us that in general simultaneous events do not appear to happen at the same time. 

A clearer example is gravitational lensing. Observed images of distant objects are distorted by the warping of space-time due to massive objects, producing appearances that differ significantly from the actual geometry of the source. Had this phenomenon been discovered prior to the advent of general relativity, it would have been highly problematic for prevailing theories. However, general relativity prepared us with a solution: appearances are shaped by space-time geometry and the observer's frame, and they cannot be assumed to directly reflect reality.

There is a clear distinction between ``what is'' and ``what appears to be''; between \textit{noumena} and \textit{phenomena}. The term ``noumenon,'' borrowed from philosophy, refers to the actual, underlying reality---the state of the universe and the real things in it, independent of the observer's frame of reference---while ``phenomenon'' refers to the appearances shaped by space-time geometry and the observer's perspective. While phenomena are what we observe, as the lessons of relative simultaneity and gravitational lensing tell us, it is important to be careful when attaching direct noumenological significance to them.

When this distinction has been ignored, the consequences have been profound. Geocentrism, for instance, conflated apparent retrograde planetary motion with actual motion (phenomena = noumena), leading to the invention of epicycles and other constructs to reconcile observations with a fundamentally flawed paradigm. When it was realised that a ``simple'' system involving uniform motion around circular epicycles travelling uniformly around circular deferents couldn't accurately account for phenomena, Hipparchus moved the Earth from the centre of the geocentric model (!) to an offset point called the eccentric. And when that model still could not accurately account for apparent planetary motion, Ptolemy modified it again by adding the equant, a separate point offset from the centre of each planet's deferent, about which a planet's epicycle would have the \textit{appearance} of uniform motion around its deferent.

This final tweak by Ptolemy provided enough accuracy that the model remained in use without any further significant change for another $\sim1400$~years. But it came at the price that the Aristotelian principle of uniform circular motion was rendered phenomenological, while the apparent retrograde motion of the planets remained noumenological. The heliocentrists in the sixteenth and seventeenth centuries would fundamentally alter this description, with a model that finally, correctly, explained that the retrograde motions of Mars, Jupiter and Saturn were actually phenomenological, not noumenological. This was Copernicus's most important insight, and it paved the way for several other key insights like Galileo's concept of inertia, a more accurate description according to Kepler's laws, and ultimately to synthesis in Newton's laws of motion and his law of universal gravitation.

The distinction between noumena and phenomena, and the importance of recognising that phenomena are not always directly tied to noumena, has been understood in principle for millennia, though in practice we usually tend to favour descriptions that draw a trivial connection between the two. This is the observation at the heart of Plato's allegory of the cave, in which shadows on a wall (phenomena) are described by chained people who cannot see the fire and actual objects that cast the shadows behind them, and therefore draw false conclusions that the shadows themselves are the essential noumena.

These examples reveal the profound limitations of assuming that appearances directly represent noumena. This assumption, while intuitive and useful as a starting point, has historically led to significant problems, necessitating paradigm shifts to distinguish phenomenological appearances from the underlying noumenological reality. When we assume that what we see accurately represents what is really happening, it leads us down wrong paths. This distinction, critical to understanding the universe's underlying structure, has often been ignored throughout history, leading to significant misinterpretations of apparent phenomena. In light of this, we would do well to formally acknowledge the counter-position, that \textit{the universe does not necessarily appear exactly as it actually is}.

In the history of astronomy, one of the most important lessons has been that apparent phenomena often mislead us into interpreting them as noumenal realities. This is the essence of the cosmological principle, learned the hard way: the core of this principle is not simply the lesson that we should not assume we occupy a special place in our universe, but that the appearance of the world from our isolated perspective has often misled us when we interpret phenomena as direct representations of its true state. The world appears geocentric. Retrograde planetary motion seems most obviously to represent noumenal motion rather than a purely phenomenological effect. The apparent distribution of stars in our galaxy initially suggested that we are at the centre of the Milky Way. The heliocentrists during the Scientific Revolution showed that the apparent retrograde motion of planets could be more accurately explained by a model that treated it as a purely phenomenological effect of Earth's motion around the Sun. Similarly, Shapley's study of globular cluster distributions disproved the interpretation that we are at the centre of the Milky Way, revealing the effects of interstellar dust and extinction and placing the galactic centre far from our location.

Later, Hubble's discovery of the redshift-distance relation presented another challenge. Initially, this observation could have been interpreted as indicating that Earth was at the centre of a cosmic explosion, with galaxies receding faster the further they were from this central point. But this geocentric interpretation was rightly rejected as inconsistent with the lessons of the past. Instead, the redshift-distance relation was understood as a result of the large-scale isotropy and homogeneity of the universe. In this interpretation, the apparent recession of galaxies is explained by the expansion of space itself: in an expanding universe, every observer witnesses the same redshift-distance relation, as the further away a galaxy is, the faster it appears to recede.

This issue is particularly glaring in a general relativistic framework, where space-time is inherently warped. Observed phenomena are frame-dependent, shaped by the geometry of space-time, and often deviate significantly from the underlying reality. General relativity fundamentally describes gravitational phenomena through the warping of space-time, decoupling appearances from noumenal structures and emphasizing their dependence on the observer's frame. General relativistic phenomena come to us like reflections in fun house mirrors: distorted, counterintuitive, and shaped by perspective rather than directly reflecting the true underlying reality. Yet the synchronicity assumption of FLRW cosmology, by tying apparent phenomena within an expanding general relativistic universe directly to noumena, fundamentally contradicts this basic feature of general relativity. In a properly relativistic cosmology, observed phenomena must be understood as emergent, frame-dependent properties shaped by space-time geometry, in which case an assumption that they are direct reflections of the noumenological structure of the universe should be recognised as an unnatural contrivance. This recognition highlights the limitations of the synchronicity assumption, pointing to the need for a cosmological framework that allows phenomena to emerge naturally from the underlying geometry of the universe.

In our original interpretations of planetary retrograde motion and the appearance that we are at the centre of the Milky Way, we swung and missed, tying phenomena directly to noumenal realities and missing the curveball reality had thrown at us. But with Hubble's law, we finally connected. We avoided tying the redshift-distance relation directly to a noumenal explosion and instead adopted a more general interpretation grounded in the cosmological principle that we do not hold a place of particular importance in the universe. Yet, while we made contact, the hit went foul. The synchronicity assumption of the FLRW model, which ties phenomena directly to noumena by requiring a synchronous universe in the cosmic rest frame, means we ultimately misjudged the curveball reality threw at us again. The lesson we failed to fully internalise is that apparent phenomena are often not directly tied to noumenal structures, and that this is particularly unexpected in a relativistic universe. With two strikes already on the count, we remain at bat, faced with the opportunity to reconsider this assumption before history hands us a third.

These lessons underscore the need for a cosmological framework that avoids conflating phenomena with noumena, allowing apparent phenomena to emerge naturally from a deeper noumenological structure. For these reasons, we reject this prior assumption and explore a cosmological model that does not carry this unnatural feature. Rather than presupposing a direct and simplistic connection between phenomena and noumena, this model constructs a framework where these connections are intricately refined, allowing for a more natural explanation of a wide variety of phenomena. By decoupling these aspects, the model provides a sophisticated and coherent approach to understanding the universe, better reflecting the principles of general relativity and the evidence of cosmic evolution.

Over the past $\sim$half-century, unresolved problems in cosmology have accumulated, often stemming from limitations in the standard FLRW framework, which cannot explain why certain features of our universe \textit{should} be as we've found. Could it be that by insisting on synchronous cosmic models, and by extension on the principle that cosmological phenomena trivially reflect our three-dimensional universe and the noumena that occur within it, we have overlooked alternatives capable of explaining why the universe is expanding, why the cosmological horizon appears uniform in all directions, why observations suggest spatial flatness, and why the universe is currently accelerating? Inflation \cite{Guth1981} was proposed to address some of these challenges---most notably, it provides an elegant solution to the flatness problem and the horizon problem, which is linked to the issue of nonlocality discussed in \S~\ref{sec:relativity}. More recently, early dark energy has been proposed as an additional modification of the standard model, to address the Hubble tension problem \cite{Kamionkowski2023}. But these are ad hoc approaches, much like Hipparchus's eccentric and Ptolemy's equant were to geocentrism: they resolve specific problems within the prevailing framework but at the cost of introducing additional layers of complexity without revisiting the foundational assumptions that create the problems in the first place. This insight suggests that by clinging to synchronous models, and by extension to the assumption that cosmological phenomena directly reflect noumenological reality, we may have overlooked the true nature of the universe---one in which apparent phenomena emerge naturally from a deeper structure. 

Allowing the possibility that the three-dimensional universe need not be synchronous in the cosmic rest frame, I have previously discussed a promising alternative \cite{Janzen2012,Janzen2015}. By discarding the synchronicity assumption, which is at odds with the principles of general relativity, and embracing the distinction between noumenological structure and phenomenological appearances, we gain a deeper perspective on the universe's fundamental nature. This cosmological model is fundamentally different from standard relativistic cosmology, addressing the unresolved problems of FLRW while providing a natural explanation for observed phenomena---such as the flat $\Lambda$CDM expansion rate---as emergent, phenomenological properties that align with observations and arise naturally from a deeper noumenological structure. It is guided by two key insights derived from the cosmological interpretation of relativity laid out above. These key insights build directly on the shortcomings of FLRW cosmology, addressing unresolved tensions in its foundational assumptions and providing a framework for understanding the universe's evolution in a manner consistent with general relativistic principles. They address two critical aspects of cosmology: the nature of gravitational collapse and the relationship between the three-dimensional universe and four-dimensional space-time. Together, they form the foundation of an alternative framework that reinterprets apparent phenomena as emergent properties of a deeper noumenological structure.

The first key insight relates to gravitational collapse. As shown in \S~\ref{sec:BH}, gravitational collapse, even in the ideal spherically symmetric case, does not reach an event horizon before infinite cosmic time has elapsed. Thus, complete collapse does not occur in our universe. However, in the Eddington-Finkelstein frame, infalling particles \textit{do} cross the horizon smoothly in finite \textit{proper} time. This inevitability along infalling particles' worldlines---the certainty that matter \textit{now} in our universe \textit{will}, at some finite proper time in its future, find itself beyond the boundary of infinite cosmic time---suggests that physical reality must extend \textit{beyond} our universe. Something \textit{will} persist \textit{after} our universe ends infinitely far into its future. Otherwise, infalling particles' worldlines would reach abrupt ends in finite proper time, at points where the relativistic space-time geometry smoothly connects to a future they are destined to exist in. The products of this process in the infinite future, however, will differ greatly from the black holes currently observed, which are dominated by matter slowly approaching their present horizons.

In the infinite future, processes like dynamical friction and radiation will cause gravitationally bound groups of galaxies to lose kinetic energy and coalesce into a single black hole at their centre of mass. For our Local Group, this final black hole will result from the mergers of the Milky Way, Andromeda, Triangulum, and the many dwarf galaxies they consume, much as the Large and Small Magellanic Clouds are being absorbed by the Milky Way today. The resulting elliptical galaxy will radiate away energy through stellar processes and dynamical interactions. Over time, neutron stars and stellar black holes will merge and spiral into the central supermassive black hole. Eventually, the energy in the Local Group will either fall into this collapsing mass or radiate away, leaving a final supermassive black hole that, at the end of cosmic time, will cross its event horizon and collapse into a future beyond our universe. 

This smoothness of the global geometry at the event horizon---the continuation of all the infalling particles' worldlines through the event horizon in finite proper time---thus implies that reality extends beyond the finite futures of infalling particles in our universe, even though they only reach the horizon at the end of cosmic time. This conclusion underscores the necessity of considering infinite cosmic time as a fundamental property of our universe. It also hints that our own universe may have originated through gravitational collapse in a prior one, after collapsing matter there crossed an event horizon.

While the first insight addresses the inevitability of gravitational collapse and its implications for cosmic time, the second insight explores the deeper relationship between the dynamically evolving three-dimensional universe and the four-dimensional space-time in which it resides. By understanding how these dimensions interrelate, we gain a clearer view of the cosmological principles that underlie the apparent coherence and consistency of our universe across cosmic scales.

This insight involves the distinction between the dynamically evolving three-dimensional universe and four-dimensional space-time in the cosmological interpretation of general relativity. The traditional view that matter warps space-time and space-time tells matter how to move must be reassessed when space-time and reality are seen as distinct aspects of the same physical theory. In this new paradigm, the cosmological description of the three-dimensional universe lays out events throughout eternity within four-dimensional space-time. Simultaneously, local space-times emerge where matter curves space, warps time, and influences particle paths. Careful consideration is needed to reconcile these two perspectives.

In the cosmological frame, the universe is \textit{always} a well-defined ``knife-edge'' within space-time. This is the noumenological reality of the universe: a three-dimensional space that evolves dynamically from one cosmic instant to the next. It is incorrect to think of it as warped \textit{cosmologically} in a way that disrupts its global structure: at every moment of cosmological time, the noumenological universe iz present; it iz not warped into its past or future, as those events zare at different cosmic times. The cosmic ``now,'' for all eternity, is a well-defined three-dimensional space. However, the noumenological universe is locally warped by matter, creating phenomena---e.g., gravitational lensing and waves---that reflect the frame-dependent distortions of this underlying structure. Local proper time, too, evolves differently from cosmic time, reflecting a \textit{layered} structure of space-time. Local general relativistic solutions describe these spatial and temporal curvatures phenomenologically, capturing the effects of local warping and time dilation as manifestations of the noumenological evolution of space in cosmic time.

These descriptions are not inconsistent. Recall the thought experiment from \S~\ref{sec:BH} involving a trampoline---a toy model with balls rolling, merging, and causing gravitational ripples. At every instant, since there zare no singularities, the noumenological universe iz diffeomorphic to a flat plane. \textit{Cosmologically}, the space in this toy model can be described as perfectly flat even while it is \textit{actually} being dynamically warped by matter. If the trampoline is infinitely large and the universe relativistic, the cosmological space-time solution must be Minkowski space, despite the noumenological warping of space at all cosmic times by matter. This diffeomorphism represents a differentiable, bijective mapping that ensures the noumenological warping of space by matter at all cosmic times does not disrupt the global equivalence to Minkowski space in the cosmological sense. This equivalence guarantees that the cosmological foliation of space-time remains well-defined, even as local space-time warping---a phenomenological effect arising from the interaction of matter with gravitational fields and described in various proper frames---produces observable phenomena such as lensing and ripples. This differentiability constraint ensures the integrity of the cosmological structure, which must remain singularity-free and hence differentiable, as justified by the physical arguments for the cosmological interpretation of gravitational collapse discussed in \S~\ref{sec:BH}.

This distinction, between cosmological space-time and space at all cosmic times, highlights a key difference between Einstein's operational approach to relativity and the cosmological interpretation presented here. Einstein's operational approach conflates space-time and physical reality as a single entity, equating the phenomenological geometry of space-time with noumenological reality. In this new framework, existing reality and space-time are formally distinct. As such, space may still \textit{really} warp through gravitational processes, but its cosmological evolution can---though not necessarily must---be described by space-times that are vacuum solutions of Einstein's equations. Einstein's operational approach, which conflates space-time with physical reality, fundamentally precludes the separation of the large-scale phenomenological space-time geometry from the more basic topology of an enduring, three-dimensional universe with a well-defined metrical connection---as illustrated by the global Minkowski space-time in the above example.

This second insight raises a critical question: must the dynamical evolution of the universe necessarily be determined by the average densities of matter and radiation filling space? If general relativity is fundamental, cosmological space-time must satisfy Einstein's equation. Local evidence strongly supports the idea that matter influences the phenomenological evolution of space. However, if existence is a fundamental property of the noumenological three-dimensional universe, with cosmological time universally well-defined across gravitational fields, voids, solar systems, and black holes, it seems implausible that this evolution depends solely on the average densities of matter and radiation. As discussed in \S~\ref{sec:relativity}, the universal consistency of cosmological time is critical for maintaining the imprinted anisotropies of the CMB over 13.8 billion years. Such coherence requires an absolute cosmic time and a globally consistent expansion rate, hinting at a non-local global structure. This introduces a philosophical tension with theories that seek to explain cosmological structure purely through local averages of matter and radiation. Rather than being determined by local densities, the absolute nature of cosmic time and expansion suggests a more fundamental structure underlying the evolution of the universe. Such a structure would imply that expansion is not simply a response to fields but an intrinsic property of the universe's geometry itself, embedded in its layered structure. 
An absolute structure of this kind would extend beyond the explanatory reach of inflationary models, suggesting the need for a deeper framework. Instead, we might explore whether the cosmological space-time geometry arises more fundamentally---through the process that created our universe, from which the emergence of fields and their associated dynamics would later arise.

This layered understanding of time, space, and expansion highlights the apparent invariance of cosmic expansion across scales. Locally, gravitationally bound structures---from molecules and stars to galaxies and clusters---do not participate in the Hubble flow. Within these regions, gravitational forces dominate, resisting cosmic expansion entirely. Beyond these thresholds, however, the Hubble flow emerges as a constant, uniform expansion rate that governs the large-scale dynamics of the universe. Observationally, the Hubble flow appears invariant across space, even between gravitationally influenced regions, such as within the Local Group or between the Local Group and the Coma Cluster. Deviations from Hubble's law observed in redshift-distance relations reflect the gravitational influence of nearby structures and resulting local motion \textit{through} space, not a variation in the expansion rate itself. If expansion were influenced by local structures, the anisotropy signature imprinted on the CMB could scarcely be thought to have survived 13.8 years of travel through space whose expansion rate not only varied in cosmic time, but also based on local density distributions.

This threshold effect is tied to a critical density gradient, below which gravitational forces hold structures together and above which the expansion dominates. For example, within the Local Group, galaxies are gravitationally bound, while at the supercluster level, the Hubble flow can be subtracted from redshift-distance relations to reveal the watershed-like structure of large-scale flows, as demonstrated by \citet{Tully2014} in their analysis of the Laniakea Supercluster. They mapped peculiar velocity flows to identify regions where galaxies' motions converge, defining superclusters as volumes within such surfaces. This method illustrates how, at supercluster scales, the intersection between local gravitational interactions and global expansion becomes evident. Local structures within these vast regions may exhibit complex motions due to gravitational interactions, yet the overall supercluster participates in the Hubble flow. By removing the mean cosmic expansion, \citet{Tully2014} revealed the underlying watershed-like structures, highlighting the layered nature of cosmic expansion.

The boundaries of gravitationally bound structures introduce a significant challenge regarding this intersection between local gravity and the Hubble flow. If local gravity halts the Hubble flow at these thresholds, what happens if a cluster subsequently evaporates? For example, if the density of an evaporating cluster drops to the mean density of the Laniakea Supercluster, would the space between its galaxies begin to expand immediately, or would intermediate dynamics emerge? This issue underscores a broader problem with the standard FLRW big bang cosmology, which describes cosmic expansion as an inertial property of the universe---a legacy of the big bang. In \textit{all} FLRW models containing matter and radiation, these universes necessarily begin by decelerating due to dominant gravitational effects, continuing to expand at later times only because they \textit{were already} expanding. This expansion problem, articulated by \citet{Eddington1933} (see \S~\ref{sec:intro}), highlights that such descriptions fail to explain why these universes \textit{should} expand. Without such a mechanism, the standard model also provides no means for expansion to restart locally once it has been halted, and it is unclear how the global Hubble flow interacts with the edges of bound structures. By contrast, the layered geometry hypothesis posits that expansion is not an inertial legacy of the big bang but an absolute, global property of the universe, inherently tied to its fundamental topology. This insight fundamentally reinterprets the nature of cosmic expansion, suggesting that the observed flat $\Lambda$CDM expansion rate emerges naturally from the layered structure, rather than requiring ad hoc modifications like inflation or early dark energy. It suggests that even as densities drop or gravitationally bound structures evaporate, expansion would persist as a universal feature, unaffected by local dynamics outside of threshold densities.

These intersections between local and global dynamics underscore the layered structure of observed expansion dynamics, exposing deeper tensions within the standard model. Beyond challenging the synchronicity assumption, this framework questions the core premise of FLRW cosmology: that the universe's expansion rate is determined primarily by fields like the average densities of matter and radiation. Observations like the CMB anisotropy signature and supercluster flow dynamics indicate that outside of bound structures, the expansion of space may be an absolute, global property of our universe. Much like cosmic time appears to be an absolute property, according to which our three-dimensional universe is well-defined within a four-dimensional foliated space-time geometry, it may be that the universe also has an absolute expansion rate that is unaffected by local gravitational interactions except where densities exceed a threshold. These observations support the hypothesis of a layered geometric structure, in which cosmic time, space, and expansion are interconnected through a basic four-dimensional topology. This foliated metric space provides the absolute structures of time, space, and expansion that underlie and govern the evolution of phenomenological dynamics in our universe. Rather than being secondary effects of local densities or motions, these absolute structures act as the framework within which local phenomena occur, ensuring coherence across cosmic scales. This hypothesis, while still in its exploratory stages, provides a coherent framework that reinterprets observed phenomena as natural consequences of a layered, absolute structure. Its strength lies in its ability to unify disparate observations under a single geometric paradigm, opening avenues for further theoretical development and empirical testing.

\subsection{Synthesis of a new cosmology from first principles}\label{sec:cosmo_synth}

Building on the insights developed in the preceding section, we now synthesize a cosmological model grounded in first principles. Central to this effort is the layered geometry hypothesis, informed by empirical evidence and the goal of developing a logically well-motivated and coherent explanation of various observations:---from the coherent description of existence afforded by the cosmological interpretation of relativity and its associated cosmic time; the scientific evidence from cosmology verifying the principle of absolute cosmic time and demonstrating that relativistic space-time is not the underlying reality but rather a map of events occurring within the evolving three-dimensional universe; the observation that the synchronicity assumption of FLRW cosmology is ill-motivated general relativistically and precludes more natural universes where phenomena are not trivially tied to noumena; the isotropy of the CMB, which supports the consistency of universal laws across scales; and the evidence that the expansion of space is absolute outside gravitationally bound structures, as shown in CMB and supercluster studies. The layered geometry hypothesis posits that the universe's large-scale dynamics emerge from the interplay between absolute structures of cosmic time, space, and expansion within a four-dimensional foliated metric space. This reinterpretation challenges conventional constructs of relativity and cosmology, emphasizing foundational principles over operational assumptions. Specifically, it: (i) decouples the existence of space-time from the distribution of matter, moving beyond the constraints of Mach's principle; (ii) suggests that the cosmological solution may be fundamentally independent of matter fields, potentially arising as a vacuum solution of the Einstein field equations; (iii) distinguishes between cosmological space-time and the real, existing three-dimensional universe by anchoring the latter in a particular foliation of the former, such that real ``space'' at any instant is diffeomorphic to slices of the cosmological geometry;\footnote{This distinction holds only if the warped ``real spaces'' are Riemannian manifolds, allowing smooth mapping to the cosmological hypersurface. The discussion of gravitational collapse in \S~\ref{sec:BH} is essential to this argument, as it establishes that singularities cannot form in a cosmologically relativistic universe.} and (iv) permits local spatial evolution in accordance with the full Einstein field equations, accommodating nonzero stress-energy densities.

Our objective is to define a cosmological foliation that serves as the foundational layer for the noumenological universe. This foliation provides the topological framework anchoring the universe's geometry---an enduring three-dimensional reality where mass influences curvature, gravitational waves propagate, and space evolves dynamically. At any cosmic instant, this real space is a Riemannian manifold with a positive-definite metric signature, geometrically consistent with and smoothly mapped onto the cosmological hypersurface of the foliated metric space. This anchoring of the noumenological universe within the cosmological foliation ensures that the coherence of physical laws and the universality of phenomena across scales are embedded within the foundational structure of the universe itself. Building on this foundation, we construct a cosmological model guided by geometric coherence and informed by heuristic principles derived from physics.

Physics indicates that the topological basis of such a universe must possess significant structure. Universal physical laws appear to act uniformly across space and time, enabling principles derived on Earth to describe phenomena at the farthest reaches of the universe. Observations of universal coherence and isotropy compel the adoption of deSitter space as the specific topological basis for cosmology, for several reasons. This perspective aligns with Eddington's argument that $\Lambda$ provides a universal scale for the radius of curvature of the vacuum, $\sqrt{3/\Lambda}$, and that material structures derive their measurable sizes from comparison with this intrinsic length scale \cite{Eddington1923}. As Eddington noted, ``An electron could never decide how large it ought to be unless there existed some length independent of itself for it to compare itself with.'' Despite the elegance of Eddington's reasoning, his interpretation of $\Lambda$ as a fundamental constant has been overshadowed by Einstein's evolving views on the cosmological term during the 1920s and his later dismissal of it as his ``greatest blunder'' \cite{Janzen2014}. This dismissal, influenced by de~Sitter's vacuum solution and Einstein's adherence to Mach's principle, has shaped the contemporary view of $\Lambda$ as an empirical parameter rather than a fundamental feature of space-time. However, as Eddington argued, $\Lambda$ provides the geometrical scaffolding necessary for the uniformity of physical laws, linking the intrinsic structure of space-time to the consistency of material phenomena across the cosmos. Observations such as a photon emitted near the dawn of time arriving intact and consistent with well-defined interactions underscore the predictive power of this framework. In the absence of such basic structure, why shouldn't such a photon arrive at Earth in the form of a nightmare or a hot bowl of soup---a metaphorical or even literal elephant in the room? Such evidence supports the conclusion that the topological basis of cosmological space-time should not only be a metric space, but one that is maximally symmetric \cite{Janzen2012}.

We construct this metric space using a real coordinate basis. Relativity implies that this space must have a Lorentzian signature. Instead of introducing this signature ad hoc, such as through a Wick rotation applied to a general maximally symmetric space, a more fundamental approach emerges. Remarkably, the only four-dimensional Riemannian manifold constructed from a purely real coordinate basis with both maximal symmetry and an intrinsic Lorentzian signature is de~Sitter space \cite{Janzen2012}.\footnote{The term ``Riemannian'' here highlights the absence of a fundamental distinction between this space and other real maximally symmetric spaces constructed from first principles. The Lorentzian signature of de~Sitter space arises naturally from the real coordinates and finite curvature constant, distinguishing it from spaces with ad hoc Lorentzian or pseudo-Riemannian signatures derived via Wick rotation.} This result aligns with cosmological evidence that our universe includes a cosmological constant, suggesting $\Lambda$ as a fundamental constant of space-time topology. Eddington's principle---his insight that the coherence of physical laws and their consistency across space and time reflects a fundamental geometric structure---provides further empirical support for this framework. Guided by these considerations, we adopt de~Sitter space as the topological basis of cosmological space-time.

Four-dimensional de~Sitter space can be visualised as a hyperboloid of one sheet embedded in five-dimensional Minkowski space, with the hyperboloid's axis aligned with the time axis of the embedding space. A hyperboloid of one sheet is a ruled surface: in three-dimensional Euclidean space, imagine tilting a bundle of straight lines 45 degrees within the plane perpendicular to a radial direction and rotating this radial line about the origin. As the tilted lines pivot and drop around the circle, they form a two-dimensional hyperboloid of one sheet. Four-dimensional de~Sitter space is similarly generated by rotating a bundle of (tilted) null lines perpendicular to the time axis of a five-dimensional Minkowski embedding space. Because the tilt direction is unrestricted, this hyperboloid is doubly ruled, with two complementary sets of perpendicular bundles of null lines covering its surface.

In the Minkowski embedding space, four-dimensional de~Sitter space is a hypersurface constructed from invariant null lines. This hypersurface describes a 3-sphere that exponentially contracts from infinite ``time'' in the past to a finite radius, then exponentially expands into the infinite future, as described by the line-element, 
\begin{equation}
    ds^2=-dT^2+\frac{3}{\Lambda}\cosh^2\left(\sqrt{\frac{\Lambda}{3}}T\right)\left(\frac{dR^2}{1-R^2}+R^2d\Omega^2\right).\label{eq:dS_hyp}
\end{equation}

\textit{Definition of cosmological line-element.} The contracting and expanding 3-sphere defined by the foliation of de~Sitter space in Eq.(\ref{eq:dS_hyp}) serves as the foundational scaffolding of the layered geometry, anchoring the universe's noumenological structure in our cosmological model. One of the two complementary bundles of null geodesics is chosen to define the universe's fundamental comoving rest frame, while photon trajectories are defined as the geodesics that remain at rest within the 3-sphere of the de~Sitter hyperboloid. From the perspective of fundamental observers who describe themselves as being at rest, this definition alters the causal structure of the space-time described by Eq.\ref{eq:dS_hyp}.

Considering that space contracts and expands with cosmic time, a line-element that respects both this new causal structure and the cosmic foliation can be written by first defining the radius of the 3-sphere as a timelike variable $r$. Recognising that the cosmic 3-sphere is always orthogonal to this timelike direction, the line-element in the two dimensions in which the fundamental observers are not moving should remain unaltered. Allowing the metric coefficients in both the $r$ direction and the third spatial direction, $t$---through which photons ``move'' along \textit{null lines}---to depend on both $r$ and $t$, the general line-element may be written, 
\begin{equation} 
    ds^2=-B(r,t)dr^2+A(r,t)dt^2+r^2d\Omega^2,\label{eq:JB} 
\end{equation} 
where $d\Omega^2$ describes the two spherically symmetric dimensions of space that contract and expand with $r$, the universe's radius \textit{and} timelike coordinate. The radius $r$, interpreted as both a timelike variable and the evolving radius of the cosmic 3-sphere, ensures consistency between the cosmic foliation and the dynamic geometry described by the line-element. This line-element must solve the vacuum Einstein equations with a positive cosmological constant, as it fundamentally describes de~Sitter space. Solving these equations constrains the functions $B(r,t)$ and $A(r,t)$, yielding the Schwarzschild-de~Sitter (SdS) geometry: 
\begin{equation} 
    \mathrm{d}s^2=-\frac{r}{\frac{\Lambda}{3}r^3+\frac{2GM}{c^2}-r}\mathrm{d}r^2+\frac{\frac{\Lambda}{3}r^3+\frac{2GM}{c^2}-r}{r}\mathrm{d}t^2+r^2\mathrm{d}\Omega^2,\label{eq:SdS_statical} 
\end{equation} 
for which $r > 0$ is \textit{timelike} when ${\Lambda}G^2M^2/c^4 > 1/9$. Under this condition, the SdS geometry is a cosmological model.

How can a universe described by such an anisotropic space-time geometry appear isotropic to fundamental observers? This question lies at the heart of the cosmological model. The key insight is that the apparent isotropy arises not from the basic cosmological geometry, but from the perspective of the fundamental observers and their motion through the universe. It is a phenomenological description, and one in which the universe at any cosmic time is clearly not synchronous in the proper rest frame since the fundamental lines are not orthogonal to the cosmic 3-sphere. There are several interesting aspects of this model to consider.

To begin, while the slices of constant cosmic time $r$ are not isotropic and homogeneous in this frame, this results from the nontrivial evolution of fundamental observers along the given bundle of geodesics rather than from an actual inhomogeneity of space. By construction, space is always a maximally symmetric structure---i.e.\ an expanding 3-sphere. Fundamentally, this means that the anisotropy is a feature of the coordinate description rather than of the physical structure of the universe. The invariance of light speed, a fundamental principle of relativity ensured by the light postulate, combined with the fact that the fundamental observers are comoving within a space that is homogeneous and isotropic, guarantees that two-dimensional images formed by light arriving ``now'' from any given epoch must appear isotropic. Consequently, the cumulative image observed at any cosmic time---as light from more distant epochs arrives---will also appear isotropic to all fundamental observers, regardless of their location. This apparent isotropy arises naturally from the properties of the space and the symmetries of the model. While SdS describes an anisotropic space-time, it is a phenomenological model layered on the underlying de~Sitter topology with its cosmological foliation, which fundamentally ensures isotropy for comoving observers.

To make this idea more intuitive, consider an analogous system where particles swirl around a 2-sphere, maintaining constant relative distances. In such a system, the speed of light remains unaffected by the particles' motion, propagating equally in all directions relative to them. This invariance implies that light travels just as quickly away from these swirling particles in the direction of their motion as it does in the opposite direction. Relativistically, this picture is equivalent to one in which the particles are at rest while ``space'' itself rotates. Importantly, it does not matter whether one chooses to describe the motion of the particles or the rotation of ``space,'' as the relative photon paths must remain the same in either case.

This concept is directly analogous to the dynamics in the cosmological model. Here, instead of particles swirling around a 2-sphere with varying speeds dependent on latitude, comoving fundamental geodesics ``swirl'' at the same speed within the circular dimension of the Hopf fibration of the 3-sphere. Despite this motion, the same conclusion applies: relativistically, the fundamental observers are not moving with respect to one another. Consequently, by the light hypothesis, they must observe the world of comoving matter as isotropic, even though the phenomenological SdS geometry describes an anisotropic space-time. This isotropy is a direct consequence of the phenomenological SdS geometry being layered upon the noumenological de~Sitter topology, which fundamentally ensures isotropic observations for comoving observers.

The apparent isotropy of the universe, as seen by fundamental observers, contrasts with the anisotropic description of the space-time and even the spatial slices of constant cosmic time, as described from their perspective. This anisotropy arises naturally from the basic topology of de~Sitter space which is foliated by expanding and contracting 3-spheres along the one-sheeted hyperboloid. In this framework, the metrical connection of the basic topology, together with the fundamental worldlines and photon lines defined in the cosmological model, leads to two distinct causal structures. In the cosmological SdS geometry, null lines correspond to photon paths mapped onto the unmoving geodesics of de~Sitter space's expanding and contracting 3-spheres. This mapping ties the causal structures of SdS space-time to the hyperboloid topology, where incoming null lines trace smaller 3-spheres in the past, and outgoing null lines traverse larger 3-spheres into the future. These differences in photon paths reflect the anisotropic light cone structure and the evolving geometry of the universe. This anisotropy mirrors the perspective of the fundamental geodesics, where the distinction between the smaller universe in the past and the larger universe into the future naturally breaks the symmetry of the light cone, distinguishing the temporal structure of the universe's past from its future.

To better understand the phenomenological implications of the model, we examine the proper reference frame of fundamental observers, where the line-element provides insight into the interplay of geometry and physical processes. The proper reference frame of fundamental observers is defined as the frame of the bundle of geodesics at rest with respect to the vanishing of the gravitational potential in the SdS geometry, which move through $t$ over the course of $r$ purely as a consequence of the underlying geometry. The line-element in this coordinate frame was found by \citet{Janzen2015}:
\begin{equation}
    ds^2=-d\tau^2+(\partial_{\chi}r)^2d\chi^2+r^2d\Omega^2,\label{eq:SdS_flat}
\end{equation}
in which
\begin{equation}
    r(\tau,\chi)=\left(\frac{6GM}{\Lambda c^2}\right)^{1/3}\sinh^{2/3}\left(\frac{3}{2}\sqrt{\frac{\Lambda}{3}}[\tau+\chi]\right)\label{eq:r_SdS}
\end{equation}
is defined on the interval $\tau+\chi>0$.\footnote{A transformation between Eqs.~((\ref{eq:dS_hyp}) and \ref{eq:SdS_flat}) can be found by setting $g_{\Omega\Omega}=g_{\Omega\Omega}$ and $g_{TT}dT^2+g_{RR}dR^2=g_{\tau\tau}d\tau^2+g_{\chi\chi}d\chi^2$; i.e.
\begin{equation}
    \sqrt{\frac{3}{\Lambda}}R\cosh\left(\sqrt{\frac{\Lambda}{3}}T\right)=+\left(\frac{6GM}{\Lambda c^2}\right)^{1/3}\sinh^{2/3}\left(\frac{3}{2}\sqrt{\frac{\Lambda}{3}}[\tau+\chi]\right),
\end{equation}
\begin{eqnarray}
    -dT^2&+&\frac{3}{\Lambda}\cosh^2\left(\sqrt{\frac{\Lambda}{3}}T\right)\frac{dR^2}{1-R^2}=-d\tau^2\\\nonumber~&+&\left[\left(\frac{6GM}{\Lambda c^2}\right)^{1/3}\sqrt{\frac{\Lambda}{3}}\frac{\cosh\left(\frac{3}{2}\sqrt{\frac{\Lambda}{3}}[\tau+\chi]\right)}{\sinh^{1/3}\left(\frac{3}{2}\sqrt{\frac{\Lambda}{3}}[\tau+\chi]\right)}\right]^2d\chi^2,
\end{eqnarray}
which are well defined on $\tau+\chi>0$ and $R<1$.}

This particular line-element, like the overall geometry, bears its own interesting features. First, $r$ does not evolve with the proper time $\tau$ of any particular fundamental worldline but with the proper cosmic time parameter $c\tilde{\tau}\equiv\tau+\chi$, which is tilted with respect to the fundamental rest frame---i.e., \textit{this universe is not synchronous in the fundamental rest frame} (which is of course clear by its design). On this same point, while it is noteworthy that slices of constant $\tau$ are Euclidean (since $\partial_{\chi}rd\chi=dr$), that is neither surprising nor of particular relevance as a cosmological ``space.'' The lines perpendicular to the $\tau$-direction are the complementary bundle of null geodesics that are perpendicular, at the ``throat'' of the de~Sitter hyperboloid, to the selected bundle taken to represent the universe's fundamental observers (see Figure~\ref{fig:dS_SdS}). These lines are of course perfectly flat, just as the lines representing the $\tau$-direction are. But from the perspective of cosmology, these lines do not represent ``space'' at any moment, as the universe evolves by definition in the positive $\tau+\chi$-direction. Rather, these flat space-like lines extend from a finite distance at $\chi=-\tau$ to $\chi=+\infty$, with synchronous events laid out throughout cosmic time, from the singularity at $r=0$ out to $r\rightarrow\infty$. These lines are the complementary bundle of worldlines, equivalent to those of our ``fundamental observers,'' which might (one could speculate) provide a rest frame in this universe for antimatter. In that case, the worldlines of these complementary ``antimatter'' particles would be fleeting in the matter rest frame, with fundamental ``antimatter'' lines occurring instantaneously in the fundamental cosmic frame of matter.

\begin{figure*}[t]
\centering
\includegraphics[width=6.5in]{dS-SdS.png}
\caption{A single fundamental world line within a cosmological model constructed from the flat null lines of de~Sitter space, whose motion is unaffected by the curvature of space-time. The causal structure of this model is redefined so that the bundle of lines parallel to the black worldline form timelike worldlines, the bundle of orthogonal lines shown here in red describe flat spacelike slices, and the blue worldlines that are at-rest in the enduring 3-sphere become null lines. An inset shows the same set of lines displayed in the proper ($\chi,\tau$)-frame, where the `throat' of de~Sitter space is the singularity $\tau+\chi=0$. In this cosmological model, spherical spatial slices of de~Sitter space, taken to describe the universe at any cosmic time, are described by diagonal lines $\tau+\chi=$constant.}
\label{fig:dS_SdS}
\end{figure*}

A second interesting feature of this line-element is that the function $r(\tilde{\tau})$ is \textit{identical} to the scale factor of the flat $\Lambda$CDM model, 
\begin{equation} 
    \frac{a(\tilde{\tau})}{a_0}=\left(\frac{3H_0^2\Omega_{m,0}}{\Lambda c^2}\right)^{1/3}\sinh^{2/3}\left(\frac{3}{2}\sqrt{\frac{\Lambda c^2}{3}}\tilde{\tau}\right).\label{eq:scalefac} 
\end{equation} 
It cannot be overstated just how remarkable this coincidence is: the universe defined by this cosmological model---the exponentially contracting and then exponentially expanding 3-sphere described by Eq.~(\ref{eq:dS_hyp})---when described in the proper rest frame of either bundle of null geodesics that generate the full de~Sitter geometry, contracts to a ``singularity'' at $r=0$, then expands as a \textit{white hole} with \textit{exactly} the apparent rate of expansion we have empirically measured \textit{our} universe to be expanding. \textit{Looking out at this universe, at the galaxies that might form and evolve in it, and at the afterglow of its initial hot, dense state, these observers would measure their universe to have precisely the same expansion rate we have measured for ours.}

This connection to the phenomenological flat $\Lambda$CDM expansion rate is particularly noteworthy when viewed in the context of the argument in \S~\ref{sec:cosmo_analysis}. There, the critique of FLRW cosmology emphasised the pitfalls of conflating noumena and phenomena, underscoring the value of models that distinguish observational effects as emergent from deeper structures. Just as the heliocentric interpretation of retrograde planetary motion revealed it to be a purely phenomenological effect rather than noumenal motion, the flat $\Lambda$CDM expansion rate in this framework emerges as a natural phenomenological property grounded in a consistent and coherent noumenological structure. By recognising this distinction, we gain not only a deeper understanding of the universe's evolution but also a framework that resolves apparent tensions between observation and theory without the need for ad hoc assumptions.

Returning to the full geometric picture, it is interesting to note that the singularity at $r=0$ in Eqs.~(\ref{eq:SdS_statical})-(\ref{eq:r_SdS}) is not a singularity of the basic topological space. Rather, $r=0$ is the moment along all fundamental worldlines when the complementary null geodesic lies in the same plane, tangent to the ``throat'' of de~Sitter space at the point each fundamental worldline crosses. At $r=0$, the universe represents a moment of finite, minimum volume, marking a smooth transition between a contracting epoch in the past and an expanding epoch in the future. This reinterpretation reframes the singularity as a geometric continuity within the broader cosmological topology. Globally, there is a smooth transition from an epoch ``before'' $r=0$ which is described by the same line-element, with the universe contracting as $r$ decreases towards $r=0$. This transition underscores the global coherence of the cosmological foliation, bridging the epochs of contraction and expansion within a unified geometric framework.

This layered interpretation of the geometry underscores the distinction between the noumenological and phenomenological descriptions of the universe. This foundational structure underpins the isotropy, coherence, and universality of physical laws, providing the scaffolding upon which the universe's fundamental dynamics are built. The reinterpretation of $\Lambda$ as a fundamental geometric parameter, rather than an empirical anomaly, reframes it as a cornerstone of the universe's intrinsic structure, aligning observational coherence with theoretical elegance. This embedding not only reflects Eddington's principle but also establishes a cohesive framework where the laws of physics arise naturally from the universe's geometric structure, rather than being imposed as external axioms. 

By contrast, the phenomenological geometry, exemplified by the cosmological SdS solution, describes the universe's large-scale dynamical evolution in the matter frame, where noumenologically ``at rest'' worldlines correspond to null paths. This geometry captures observable phenomena, such as the apparent flat $\Lambda$CDM expansion rate and local gravitational effects, while preserving the anisotropic properties of the cosmological frame. However, the universe is not synchronous in the fundamental rest frame, reflecting the layered interplay between the noumenological structure and the causal properties of the SdS geometry. These phenomenological features emerge as a specific manifestation of the noumenological foundation, describing the evolution of the universe in a manner that aligns its global structure with the dynamics inferred from the matter frame.

This distinction extends beyond the cosmological scale to include local space-time geometries, which emerge dynamically from the same fundamental framework. Local space-time reflects the operational consequences of matter's influence on curvature, warping the geometry in ways consistent with Einstein's equations. The coexistence of noumenological, phenomenological, and local geometries underscores a profound conceptual shift: cosmological relativity separates space-time from reality itself, reconciling local and global phenomena within a unified framework. This conceptual separation not only reconciles local and global phenomena but also invites a rethinking of foundational principles in physics, paving the way for new theoretical and empirical insights.

\subsection{Outline of the cosmological theory's basic predictive and explanatory elements, with empirical tests}

The layered geometrical foundation of the cosmological theory extends naturally to broader implications and challenges in cosmology. Central to this framework is a conjecture connecting gravitational collapse to the origin of the universe. The smooth transition through $r=0$, as described by the cosmological SdS geometry, suggests a continuity between collapsing structures in the present universe and the genesis of our universe through analogous processes. This connection reframes our understanding of cosmic origins, grounding the model in a framework that addresses both dynamical evolution and the foundational properties of the universe.

Building on this foundational insight, the model provides identical resolutions to several persistent issues in modern cosmology. These include the nature of cosmic expansion, the apparent need for dark matter- and dark energy-driven dynamics, the flatness and horizon problems, and the cosmological constant problem. Each resolution emerges naturally from the layered interplay between the noumenological and phenomenological geometries, offering solutions without relying on ad hoc mechanisms or speculative fields.

In addition, the model addresses an emerging challenge in structure formation astrophysics, which we term the ``threshold problem.'' This issue arises from the sharp division between gravitationally bound regions and the global Hubble flow in the standard paradigm, a discontinuity at odds with expectations of smooth dynamical transitions. The layered geometry hypothesis resolves this tension by embedding cosmic expansion as a fundamental geometric property, intrinsically tied to the universe's noumenological structure. This perspective eliminates the need for artificial boundaries while offering a conceptual foundation for reinterpreting phenomena traditionally attributed to dark matter.

The model also suggests exploratory hypotheses that extend its foundational framework to address unresolved questions in physics and cosmology. These include a hypothesis linking dark matter to interactions between noumenological underdense regions and their phenomenological manifestations, as well as ideas regarding matter/antimatter duality and asymmetry. Potential connections to particle physics symmetries, including $U(1)$, $SU(2)$, and $SU(3)$, further highlight the model's capacity to integrate insights across disciplines, offering novel perspectives on the fundamental structure of physical reality. These hypotheses, while speculative, provide a foundation for addressing deeper questions about the interplay between the universe's geometry and its material and dynamical properties.

The testability of this framework is critical to its scientific merit. Empirical evidence remains the supreme arbiter of any theoretical framework, and the strength of this model lies in its capacity to provide falsifiable predictions. Existing data already align with its predictions, such as the mass of the observable universe exceeding the threshold required for the cosmological SdS geometry (discussed below) and a phenomenological flat $\Lambda$CDM expansion rate that is independent of material content. These alignments address pressing cosmological tensions, including the Hubble constant discrepancy and challenges in explaining early structure formation. The model also provides a novel perspective on CMB anisotropy, reframing features traditionally attributed to inflation in terms of the noumenological and phenomenological structures outlined earlier. Further empirical testing---including analyses of the sound horizon, local expansion rates, and dark matter distributions---offers promising opportunities to validate the model and refine its framework.

Scientific inquiry is traditionally associated with deductive reasoning: the generation and testing of hypotheses within an established framework. However, its logical counterpart---reduction---plays an equally essential but often underappreciated role. Reduction involves stripping away paradigmatic assumptions to isolate empirical evidence and reinterpret foundational principles. While deduction explores consequences from given premises, reduction reexamines the premises themselves, providing the groundwork for transformative paradigm shifts. Historically, this approach has been the foundation of major scientific revolutions, from Copernicus's heliocentric model to Einstein's general relativity.

The reductive approach guided this paper's exploration of relativity, existence, and cosmology. By considering the operational definition of simultaneity and associated interpretation of relativity in light of empirical evidence from cosmology and an associated description of physical existence that agrees with tacit experience, we reduced this interpretation to reveal its weak empirical support and its inconsistency with a broader cosmological framework. This analysis supported the cosmological interpretation of relativity, with its absolute simultaneity, as a more objective and coherent explanation of the empirical facts. This approach exemplifies the power of reduction to disentangle paradigmatic assumptions from empirical facts, paving the way for foundational insights into the layered geometry of the universe.

A similar reductive approach applies to the inflationary paradigm, separating empirical evidence from its explanatory framework. Observations of coherence, adiabatic expansion, and near scale-invariance in the CMB spectrum are often cited as direct evidence of inflation. However, these features are empirical facts, not exclusive to the inflationary explanation. Reduction reveals that these observations require a theoretical framework capable of producing exponential expansion and coherent perturbations but do not inherently demand the inflationary mechanism itself. By distilling the raw empirical data from the paradigm's interpretation, reduction opens the possibility for alternative explanations. In the model presented in \S~\ref{sec:cosmo_synth}, these features should similarly arise from the noumenological exponential expansion in the vacuum frame, Eq.~(\ref{eq:dS_hyp}), and the phenomenological apparent flat $\Lambda$CDM expansion rate in the matter frame, Eq.~(\ref{eq:r_SdS}), illustrating how a fundamentally different paradigm can address the same evidence while resolving tensions inherent to the inflationary framework.

Reduction not only separates empirical facts from paradigm assumptions but also provides a pathway for paradigm shifts. Einstein's reflection on the scientific process highlights the importance of bridging the gap between foundational principles and empirical evidence: ``The observed fact is undoubtedly the supreme arbiter; but it cannot pronounce sentence until the wide chasm separating the axioms from their verifiable consequences has been bridged by much intense, hard thinking'' \cite{Einstein1954}. This bridging process reflects the essence of reduction, requiring scientists to reexamine foundational principles and rigorously connect them to observations. By enabling the reinterpretation of observational evidence in ways that reveal deeper, unifying principles, reduction serves as the foundation for transformative shifts in scientific paradigms.

This paper's reductive methodology exemplifies this process, challenging entrenched assumptions and offering a layered geometrical model that not only resolves existing tensions, such as the identical resolutions to several long-standing cosmological challenges mentioned above, but also expands the horizons of theoretical and experimental physics. By embracing reduction alongside deduction, the framework provides a rigorous basis for rethinking the universe's foundational properties.

The remainder of the section begins with a detailed analysis of the reductive methodology applied to the layered geometry hypothesis, providing a worked example of how this process identifies tensions, separates empirical evidence from paradigmatic assumptions, and leads to new principles. It then explores a conjecture connecting gravitational collapse to cosmic origins, discusses identical resolutions to major cosmological problems, introduces exploratory hypotheses, and concludes with an examination of empirical testing and its critical role in evaluating the model.

\paragraph{Reduction in cosmology: logical foundations of the layered geometry hypothesis.} 
Reduction disentangles empirical evidence from the interpretive frameworks of existing paradigms, enabling foundational insights and fostering paradigm shifts. In the context of cosmology, this process reexamines observational data---separating empirical facts from their paradigmatic interpretations---and constructs alternative principles that unify and extend these observations under a coherent framework. A key example explored in this paper is the reductive analysis of the flat $\Lambda$CDM expansion rate, which provides critical insights into the limitations of the FLRW paradigm and motivates the layered geometry hypothesis.

The standard FLRW paradigm interprets cosmic expansion as driven by energy densities dominated by dark matter and dark energy at different cosmic epochs. Within this framework, the flat $\Lambda$CDM expansion rate reflects the interplay of these ``dark'' fields, invoked to account for phenomena that appear inconsistent with the standard matter-energy content of the universe. However, reductive analysis reveals that this interpretation conflates empirical observations with their theoretical explanations. The observational evidence itself is remarkably specific: the universe's expansion history is consistent with an initially decelerating and subsequently accelerating rate, precisely described by the flat $\Lambda$CDM scale factor, Eq.~(\ref{eq:scalefac}).

Moreover, this expansion rate appears to be absolute beyond the boundaries of gravitationally bound structures, as evidenced by data from the CMB, superclusters, and local observations. The empirical evidence further suggests a threshold where local gravitational forces intersect with the global expansion rate, indicating a nonlocal structure underpinning the dynamics of apparent expansion. These observations, when disentangled from their interpretation within the FLRW framework, provide a clearer picture of the phenomena that any cosmological model must address.

This reductive insight exposes significant tensions within the FLRW paradigm, encompassing several interrelated issues:
\begin{itemize}
    \item[i.] \textit{Reliance on local mechanisms for global coherence}: The paradigm's reliance on local, field-driven mechanisms to describe expansion is fundamentally at odds with the empirical evidence supporting global coherence. This foundational incompatibility motivates the introduction of the inflation scenario to address nonlocal coherence issues, such as the horizon problem. However, the reliance on inflation underscores the paradigm's deeper misalignment with observations.
    \item[ii.] \textit{Dependence on ``dark'' fields}: The paradigm's correspondence between energy densities and expansion rate necessitates the invocation of ``dark'' fields to explain expansion. While dark matter has some independent observational support, dark energy remains an ad hoc construct with unresolved issues, including its unknown physical nature and lack of direct evidence.
    \item[iii.] \textit{Intrinsic challenges with cosmic expansion}: The FLRW model treats cosmic expansion as an inertial legacy of the big bang, which fails to explain why the universe should expand at all. Problems such as the flatness and cosmological constant challenges remain unreconciled within this framework.
    \item[iv.] \textit{Synchronicity assumption}: The paradigm's synchronicity assumption, motivated by operational relativity rather than cosmological general relativity, imposes a direct correspondence between phenomena and noumena. This assumption contradicts both historical lessons, which caution against conflating appearances with reality, and general relativity, which suggests that such direct ties are unnatural in a relativistic universe.
    \item[v.] \textit{Cosmological GR insights}: General relativity, in its cosmological interpretation, decouples space-time from reality. This interpretation recognises that the noumenal universe need not correspond to hypersurfaces of a coherent cosmological space-time metric. This decoupling, combined with other tensions, motivates the exploration of alternative frameworks.
\end{itemize}

Recognising these tensions through reduction not only challenges the empirical adequacy of the FLRW paradigm but also motivates the exploration of alternative frameworks. By disentangling the flat $\Lambda$CDM expansion rate from its standard interpretation and examining it as an independent empirical fact, we open the door to reinterpreting it through a different set of principles.

The layered geometry hypothesis emerges naturally from this reductive process. Observations of universal isotropy, the coherence of both cosmic time and cosmic expansion, and specifically the flat $\Lambda$CDM expansion rate, when considered holistically, suggest that these features may reflect intrinsic properties of the universe's geometry rather than consequences of local energy densities. By treating the flat $\Lambda$CDM expansion rate as a phenomenological property grounded in a foundational noumenological structure, the hypothesis shifts the focus from field-driven mechanisms to the geometry of the universe itself.

This reinterpretation is grounded in the universality of physics as a cosmological principle. The hypothesis posits that the layered interplay between noumenological geometry and phenomenological effects can naturally account for the observed expansion rate, its coherence across cosmic scales, and the threshold between local and global dynamics. For example, the observed expansion rate is not energy-field driven; it emerges phenomenologically while not being a noumenological property of the universe.

This example illustrates two key features of logical reduction. First, it is inherently holistic, moving beyond the separation of raw empirical evidence from paradigmatic interpretations to contextualise observations within broader principles. In this case, the reductive process not only disentangles the flat $\Lambda$CDM expansion rate from its interpretation within the FLRW framework but also synthesises evidence of isotropy, coherence, and the nonlocal nature of expansion into guiding principles. Unlike incremental approaches that address specific anomalies within a paradigm, holistic reduction identifies patterns across disparate observations to construct a cohesive framework. By challenging deeply entrenched paradigmatic assumptions, reduction uncovers alternative principles that better align with empirical evidence and logical coherence, offering a more robust foundation for scientific inquiry.

Second, reduction contrasts sharply with the ad hoc deductive approaches that dominate modern physics. These approaches often develop hypotheses within existing paradigms to address specific anomalies, deducing consequences that fit the data they are designed to explain. While such hypotheses can extend paradigms, they frequently add descriptive flexibility without resolving underlying tensions, reinforcing the paradigm rather than questioning its validity. Geocentrism offers a historical example: successive refinements such as epicycles, eccentrics, and equants resolved discrepancies but ultimately entrenched a flawed framework. Similarly, the FLRW model's incremental modifications to address tensions like the Hubble constant discrepancy mirror this pattern. These adjustments expand the paradigm's scope while obscuring core issues, allowing basic assumptions to remain unquestioned and, over time, to harden into deeply entrenched beliefs reinforced by the model's apparent successes.

This reluctance to engage in true reduction reflects a deeply ingrained operationalist attitude within the modern physics community, where discussions of ontology and foundational principles are often dismissed as ``too philosophical.''\footnote{\citet{Einstein1954} addressed this tendency to dismiss foundational inquiries directly. He argued that theorists who undertake this work should not be ``carped at as `fanciful','' emphasising that rigorous analysis of foundational principles is far from ``idle daydreaming.'' Rather, it is an essential process of bridging the ``wide chasm separating the axioms from their verifiable consequences'' through ``much intense, hard thinking.'' His ``plea'' highlights the importance of giving theorists the intellectual freedom to reexamine the logical underpinnings of prevailing paradigms, countering the dismissive attitudes that too often obstruct transformative scientific progress.} This perspective, solidified during the successes of quantum mechanics and the standard model in the twentieth century, prioritises calculative precision over conceptual analysis, reinforcing a ``shut up and calculate'' approach to science \cite{Mermin2004}. While effective in many domains, this mindset risks constraining the scope of scientific inquiry, particularly in areas where foundational issues remain unresolved. By avoiding foundational scrutiny, operationalism enables the development of models that can flexibly fit data but at the cost of overlooking deeper inconsistencies. Reduction, by contrast, seeks to reexamine the foundational principles underpinning empirical models, inviting scrutiny of the assumptions that shape our interpretation of physical phenomena. By exposing and disentangling these assumptions, reduction provides a pathway for overcoming entrenched limitations and enabling the construction of new paradigms capable of addressing unresolved problems holistically.

The layered geometry hypothesis exemplifies the power of reduction to bridge disparate observations and principles within a unified framework. By identifying tensions in FLRW cosmology---such as its synchronicity assumption, reliance on local mechanisms for global coherence, and interpretation of expansion as energy-field driven---reduction reinterprets these features through the lens of the layered interplay between noumenological and phenomenological structures. This reinterpretation treats the flat $\Lambda$CDM expansion rate as a phenomenological property of a foundational noumenological structure, providing a logically coherent and empirically motivated alternative to the FLRW paradigm. This hypothesis lays the groundwork for further theoretical and observational exploration, including a profound reinterpretation of cosmic origins. By integrating the layered geometry into a cosmologically relativistic framework, the model provides new insights into gravitational collapse and its potential connection to the genesis of universes. Central to this reinterpretation is a conjecture linking collapsing structures to universal birth.

\paragraph{Conjectured connection between gravitational collapse and cosmic origins.}
Central to the model is a conjecture that connects gravitational collapse to the origin of the universe, offering a profound reinterpretation of cosmic beginnings. This conjecture is grounded in two complementary perspectives: what can be deduced about collapsing structures within the universe and what can be inferred from the cosmological model describing universal birth.

These perspectives converge within the framework of cosmological relativity, where the principles governing collapse also inform the properties of universal birth. On the one hand, in the cosmological SdS geometry, the smooth transition through $r=0$---where the noumenological geometry ensures continuity across the cosmological foliation---suggests a direct relationship between collapsing structures in our present universe and the potential genesis of our universe within a prior one. The noumenological framework indicates that such a transition occurs when a collapsing object forms a 3-sphere in which matter moves along flat null lines of a de~Sitter hyperboloid, converging cosmologically to a finite radius. This collapsing 3-sphere naturally transitions into an expanding cosmological white hole, providing an elegant geometrical framework for understanding the birth of a new universe.

On the other hand, we can infer key properties of the collapsing object at the end of cosmic time based on the principles of cosmological relativity. A black hole at this stage would contain all the matter and energy of gravitationally bound structures within its horizon, distributed throughout a finite spatial volume. This distribution would include an overall angular momentum, reflecting the dynamics of galaxy mergers and large-scale structure formation. Importantly, cosmological relativity ensures two critical features of this collapse: (i) the matter cannot evolve into a timelike radial distribution beyond the horizon (i.e., so that matter would ``be'' in the past and future of other particles, in a meta-temporal sense, distributed throughout the timelike radial and spatial dimensions of a ``collapsing'' space-time distribution), as this would violate causality and the foundational principles of existence in this model, and (ii) the external space of the universe cannot exist after the horizon forms, as cosmic time ends, and there is no longer a ``present'' external frame for such space to inhabit.

These insights form the basis for the conjecture: beyond the horizon, the universe transitions into a new phase described as a 3-sphere embedded in de~Sitter space. In this phase, matter continues to move along null lines of the de~Sitter geometry, while photons remain stationary within the 3-sphere. Initially, the 3-sphere collapses to a finite radius, after which it undergoes exponential expansion, mirroring the contracting and expanding topology of de~Sitter space. Thus, what exists beyond the horizon must be distributed through a finite spatial volume that continues to exist in a cosmologically relativistic sense. The collapsing volume contracts over cosmic time, with matter naturally inheriting angular momentum from the collapsing object before the horizon forms, retaining this property throughout the transition into the expanding phase of the universe. These features provide a seamless geometric and physical continuity, bridging the dynamics of collapse and universal birth in a manner consistent with the principles of cosmological relativity.

This unified framework situates gravitational collapse and universal birth within the same geometric principles that govern the model. By embedding these processes in the noumenological de~Sitter topology, the model offers a cohesive explanation of cosmic origins, linking the dynamics of the collapsing object to the properties of the new universe. This conjecture not only reframes our understanding of cosmic origins but also addresses persistent cosmological challenges. Building on this foundation, the model provides identically resolved solutions to several long-standing problems, each emerging naturally from the layered interplay between noumenological and phenomenological geometries.

\paragraph{Theoretical resolutions of long-standing problems as identities of the theorised universe.}
The conjecture of gravitational collapse and universal birth provides a natural foundation for addressing long-standing challenges in cosmology. The nature of cosmic expansion is directly tied to the transition described in the conjecture: beyond the horizon, the collapsing 3-sphere reaches a finite radius before transitioning into an expanding phase, mirroring the exponential expansion of de~Sitter space. In the phenomenological matter frame, this expansion is perceived as the flat $\Lambda$CDM rate, consistent with observations. However, this apparent expansion rate is not driven by significant dark matter or dark energy components of the universe's present energy density, as traditionally assumed. Instead, it emerges as a fundamental geometric property of the noumenological structure. By embedding these dynamics within the intrinsic geometry of the cosmological foliation, the model provides a coherent framework that eliminates the need for additional speculative fields or ad hoc parameters.

The horizon problem is similarly reframed within this context. The conjecture naturally addresses the global coherence of the universe's large-scale structure, which is often attributed to inflation in standard cosmology. Here, the noumenological geometry ensures isotropy and coherence across the entire universe, as the initial conditions of the expanding 3-sphere are inherited directly from the smooth, finite-radius state at $r=0$. The global causal structure, shaped by the noumenological de~Sitter topology, inherently provides the conditions necessary for the observed uniformity of the CMB. This embedded coherence obviates the need for inflationary mechanisms to enforce causal connection across the horizon, reframing coherence as a geometric property of the cosmological foliation.

The flatness problem is resolved phenomenologically. The universe's apparent flatness emerges naturally in the matter frame as a consequence of the flat $\Lambda$CDM expansion rate. This rate reflects the phenomenological description of the noumenological 3-sphere's evolution, requiring no fine-tuning of initial conditions. In this framework, the observed large-scale flatness is not a coincidence but a direct consequence of the layered structure of the model, which connects the phenomenological appearance to the noumenological foundation.

The expansion problem, a less commonly acknowledged challenge, is also resolved through the intrinsic properties of the model. Eddington critiqued early cosmological models, including those by Einstein and de Sitter, for their failure to account for the observed expansion of the universe, introducing the idea of cosmical repulsion as a driving force for expansion. However, this concept, like standard FLRW cosmology, relies on external mechanisms and does not fully address why the universe should expand initially rather than collapse under the influence of matter and radiation. In all FLRW big bang models, the acceleration Friedmann equation predicts significant deceleration in the presence of matter or radiation, with the deceleration terms growing exponentially as the scale factor approaches zero. This inherent issue, highlighted by Hoyle in coining the term ``big bang,'' remains unresolved within the FLRW paradigm. In contrast, the model presented here posits that expansion is a fundamental property of the universe's geometric structure, rooted in the de~Sitter topology and cosmological foliation. This perspective not only resolves the expansion problem identified by Eddington but provides a foundational explanation for the universe's initial expansion without invoking speculative fields or repulsive forces.

Finally, the cosmological constant problem is addressed at a foundational level. The model distinguishes between the cosmological constant's role in the noumenological geometry and the phenomenological expansion observed in the matter frame. While the cosmological constant determines the inherent structure of de~Sitter space, the vacuum energy it describes is not associated with the phenomenological expansion rate, which is instead a geometric property of the cosmological foliation. By reinterpreting the cosmological constant's role within this layered framework, the model resolves the discrepancy between observed and predicted values, aligning theoretical and observational perspectives.

These identically resolved solutions highlight the explanatory power of the cosmological model, each emerging as a natural consequence of its underlying structure. While these resolutions address well-established problems in cosmology, the layered geometry hypothesis also offers a conceptual framework for confronting emerging challenges. Among these is the threshold problem, an underexplored tension within the FLRW paradigm, which reflects a dichotomy between global expansion and local gravitational dynamics. This problem, whose significance is increasingly recognized, provides an opportunity to assess the broader implications of the model and explore its capacity for resolving tensions at the intersection of large-scale cosmic expansion and local structure formation.

\paragraph{The threshold problem in FLRW cosmology: conceptual tensions and resolutions.}
Modern cosmology has long grappled with the challenge of reconciling the global expansion of the universe with the local dynamics of gravitationally bound structures. Within the FLRW paradigm, both phenomena are governed by the same gravitational framework, derived from the Einstein field equations. While this connection has guided cosmological models for decades, it also gives rise to significant theoretical and empirical tensions. Chief among these is the sharp division between bound and expanding regions, suggesting two fixed and distinct regimes in stark contrast to the gradual interaction expected from gravitational dynamics. The layered geometry hypothesis, which resolves other longstanding challenges, provides a conceptual framework for addressing this problem.

The FLRW framework models the universe's large-scale dynamics by assuming isotropy and homogeneity at large scales, supported by the observed isotropy of the CMB and large-scale structure studies. The critical density serves as a key empirical benchmark, distinguishing gravitationally bound systems from regions participating in the Hubble flow \cite{Nagamine2003}. Historically, in Einstein–de~Sitter cosmology (with no cosmological constant), the critical density determined both global expansion and local gravitational binding. In this model, a universe exceeding the critical density would eventually recollapse, while one below it would expand indefinitely. In $\Lambda$CDM cosmology, where long-term acceleration is driven by \(\Lambda\), this distinction disappears, and the critical density's role is reduced to defining the universe's curvature. Despite this shift, the critical density remains conceptually linked to structure formation, where \(\Lambda\) is negligible and local binding processes dominate.

At smaller, inhomogeneous scales, FLRW predicts a nuanced interplay between local density variations and cosmic expansion. Regions with higher density should decelerate expansion more strongly, while lower-density regions participate more fully in the Hubble flow---these differences supposedly averaging out on the large scale. However, as \citet{Paraskevas2024} emphasize, the transition between bound and expanding regions lacks gradual variation, instead presenting as a stark threshold. This discontinuity contradicts expectations of continuous dynamics near the critical density, where edge galaxies should exhibit gradual transitions into or out of the Hubble flow \cite{Wang2023}.

This tension deepens at later times as the universe's average density drops below $\Lambda$. \citeauthor{Eddington1933}'s~\citeyearpar{Eddington1933} cosmical repulsion, driven by the cosmological constant, offers a potential mechanism for accelerating underdense regions, beginning in cosmic voids and gradually influencing areas closer to the critical density. Such dynamics, while speculative within FLRW, provide an essential context for considering alternative frameworks like the layered geometry hypothesis.

A deeper tension emerges when considering large-scale anisotropies in the observable universe. Observations of structures like the Shapley Supercluster \cite{Valade2024} reveal significant inhomogeneities on scales of at least $3.5\%$ of the observable universe.\footnote{The Sloan Great Wall's basin of attraction, for example, highlights significant inhomogeneity on scales of at least $500~\text{Mpc}$. Compared to the present comoving radius of the observable universe, $14.3~\text{Gpc}$, this represents $\sim 3.5\%$ of the total comoving volume.} Despite this, the global expansion rate remains perfectly isotropic and homogeneous, as inferred from CMB studies and large-scale cosmic flows. If local inhomogeneities influenced expansion rates, their effects would convolve with and overwhelm the primordial CMB anisotropy signal, which is detected at a precision of $\sim 10^{-5}$. The fact that the CMB anisotropy remains observable implies that expansion is an absolute, global property of the universe's noumenological structure, unaffected by local density variations outside bound regions.

This contradiction compounds another significant observation: the sharp transition between bound and expanding regions. Empirically, regions slightly below the critical density, such as the Local Group–Virgo region, expand at the same Hubble flow as much less dense regions, such as voids, with no gradual variation. Conversely, regions slightly above the critical density, such as the Local Group, are fully bound, showing no evidence of edge galaxies transitioning into expansion. This sharp decoupling between gravitationally bound systems and the Hubble flow \cite{Paraskevas2024, Wang2023} is inconsistent with FLRW's reliance on gravitational potential fields to describe continuous dynamics near the critical density.

The layered geometry hypothesis resolves these dual contradictions by embedding expansion as a fundamental property of the universe's noumenological structure, anchored to the de~Sitter topology. In this framework:
\begin{itemize}
    \item Global expansion persists independently of local density variations, resolving the isotropy-heterogeneity contradiction.
    \item The threshold separating bound and expanding regions emerges naturally from the layered interplay between noumenological and phenomenological structures. Bound regions dynamically resist expansion as a geometric response, eliminating the need for calibrated boundaries.
    \item As matter density drops below a critical level, cosmical repulsion in underdense regions enhances global expansion, while overdense regions experience added inward gravitational effects. These dynamics offer a novel explanation for observations traditionally attributed to dark matter, such as enhanced gravitational effects in galaxies and clusters.
\end{itemize}

Moreover, this perspective aligns with the resolution of the expansion problem as outlined in the previous subsubsection. The layered geometry hypothesis reframes expansion as an intrinsic geometric property of the universe, rooted in its noumenological structure. This provides a coherent explanation for why the universe expands at all, while simultaneously addressing the threshold problem. By unifying global and local dynamics within a single geometric framework, the hypothesis offers a resolution consistent with both theoretical principles and observational evidence.

\paragraph{Theoretical speculations and fundamental connections.} 
Building on the conceptual resolutions of emerging problems like the threshold problem, the layered geometry hypothesis invites further exploration of its broader implications. At the heart of this model is the embedding of conservation laws and the universality of physics directly within the noumenological geometry, in accordance with Eddington's principle. By rooting these foundational principles in the universe's geometric structure, the model offers a fertile framework for unifying cosmological phenomena with fundamental physics.

Beyond these grounded insights, the model invites more speculative considerations. The dual causal structures of the noumenological geometry, represented by complementary null lines, suggest a potential basis for understanding matter/antimatter duality. If matter is associated with one set of null lines and antimatter with the complementary set, this framework could provide a geometric explanation for the symmetry observed between these fundamental constituents of the universe. Additionally, asymmetry in matter and antimatter may arise naturally from the dynamics of gravitational collapse. Specifically, the angular momentum of the collapsing object prior to horizon formation might imprint a preferred handedness or asymmetry in the emergent universe, with ``matter'' inheriting the angular momentum of the collapsing structure.

The layered geometrical interpretation also hints at potential connections to the symmetries of the standard model of particle physics. The noumenological 3-sphere's structure and the properties of its foliation might offer a framework for exploring connections to $U(1)$, $SU(2)$, and $SU(3)$ symmetries, which govern the electromagnetic, weak, and strong interactions. While speculative, these potential connections open the door to a unified understanding of particle physics and cosmology through the geometry of the universe.

Taken together, these insights and speculations highlight the richness of the cosmological model, which not only resolves long-standing challenges in modern cosmology but also suggests pathways for integrating foundational physics with cosmological structure. By embedding conservation laws and the universality of physics within the noumenological geometry, the model provides a cohesive framework that bridges the gap between the fundamental laws of physics and their manifestation in the observable universe. However, as with any theoretical framework, the ultimate test lies in its ability to align with and predict observable phenomena. Thus, the layered geometry hypothesis invites rigorous empirical scrutiny, where measurable phenomena can either validate or refine its foundational principles.

\paragraph{Empirical tests of the cosmological model.} A crucial step in evaluating the validity of this cosmological model is comparing its predictions to existing empirical data. Several tests stand out, each providing opportunities to scrutinize the framework against measurable phenomena.

The first empirical test involves the mass constraint for the cosmological SdS geometry. For this geometry to qualify as cosmological, it must satisfy $\Lambda G^2 M^2 / c^4 > 1/9$, implying:
\begin{equation}
M > \frac{c^2}{3\sqrt{\Lambda G}} \approx 4.3 \times 10^{52}~\text{kg}.
\end{equation}
The calculated baryonic mass of the observable universe, $M_{b,\text{horizon},0} = 1.5 \times 10^{53}~\text{kg}$, and total mass, $M_{m,\text{horizon},0} = 9.6 \times 10^{53}~\text{kg}$, both exceed this threshold. This consistency demonstrates that the model aligns with known mass-energy constraints and provides a novel perspective for interpreting the universe's total mass distribution in light of the SdS geometry.

The layered geometry hypothesis offers a novel reinterpretation of dark matter, attributing its apparent effects to cosmical repulsion arising from underdense regions within the noumenological structure. In this framework, phenomena such as galaxy rotation curves, gravitational lensing, and large-scale structure formation emerge naturally from the interplay between local curvature and global expansion. For example, galaxy rotation curves---traditionally explained by dark matter halos---may instead reflect cosmical repulsion interacting with local gravitational potentials \cite{Rubin1980}. Similarly, gravitational lensing profiles, which measure the bending of light near massive objects, offer a means of evaluating whether cosmical repulsion can reproduce observed lensing effects \cite{Refregier2003}. On larger scales, clustering statistics and void distributions in galaxy surveys provide additional opportunities to explore this hypothesis. Distinct signatures in the observed distribution of galaxies and voids would allow researchers to distinguish between the layered geometry framework and the standard $\Lambda$CDM model.

High-precision observational campaigns, such as those conducted by Euclid and LSST \cite{Laureijs2012, Abell2009}, provide the resolution needed to rigorously test these predictions. By analyzing galaxy rotation curves, lensing profiles, and large-scale clustering within this framework, researchers can refine or falsify the hypothesis. These tests offer a pathway to integrate the hypothesis into the broader cosmological paradigm or identify its limitations, providing critical insights into the role of cosmical repulsion in shaping the universe's structure.

Structure formation offers another promising domain for empirical testing. In the layered geometry framework, removing radiation's influence on early expansion fundamentally alters the universe's initial deceleration. Without radiation dominating early dynamics, the plasma phase preceding recombination is prolonged, giving density perturbations additional time to grow under gravitational instabilities. This enhanced growth naturally accounts for the observed large-scale structures without requiring fine-tuned initial conditions or exotic mechanisms.

Simulations can model the growth of density perturbations by incorporating the phenomenological flat $\Lambda$CDM expansion rate observed in the matter frame while accounting for the noumenological evolution of comoving matter in the universe's rest frame. This layered approach ensures alignment with the hypothesis' distinct causal and geometric structures. Observational datasets, such as those from Euclid \cite{Laureijs2011} and DESI \cite{Aghamousa2016}, provide high-resolution measurements of galaxy clustering and large-scale flows that are essential for evaluating these predictions. By revisiting simulations under the layered framework and comparing them to these datasets, researchers can assess whether the model offers a more consistent explanation for early structure formation and large-scale structures.

A central feature of the layered geometry hypothesis is the distinction between noumenological evolution and phenomenological observations. Noumenologically, the universe evolves as an exponentially expanding 3-sphere embedded in de~Sitter space, with matter comoving within this structure. Phenomenologically, in the matter rest frame, this evolution appears as a flat $\Lambda$CDM expansion rate. Connecting these two perspectives is crucial for understanding the model's predictions, as mapping noumenological evolution through cosmic time to phenomenological expansion in proper time reveals insights into structure formation and cosmic dynamics.

Recalibration of effective density parameters is an essential step in testing the layered geometry hypothesis. In this framework, radiation density ($\Omega_r$) does not influence global cosmic expansion but remains critical for local and phenomenological processes, such as the evolution of the photon-baryon fluid and the formation of the sound horizon. This departure from the standard $\Lambda$CDM framework necessitates recalibrating $\Omega_r$, $\Omega_m$, and $\Omega_\Lambda$ to align observational analyses with the model's distinct causal and geometric structure.

These recalibrations are particularly important for addressing persistent tensions in cosmological data. Removing radiation's influence on global expansion delays recombination and shifts the sound horizon's evolution, potentially affecting the angular power spectrum of the CMB and the scale of baryon acoustic oscillations (BAO). High-precision datasets, such as those from Planck \cite{Aghanim2020}, DESI \cite{Aghamousa2016}, and Euclid \cite{Laureijs2011}, provide critical benchmarks for evaluating these recalibrations. Statistical methodologies, including the Bayesian Information Criterion (BIC) \cite{Schwarz1978} and Akaike Information Criterion (AIC) \cite{Akaike1974}, ensure rigorous comparisons between recalibrated parameters and standard $\Lambda$CDM predictions, offering robust tests of the hypothesis.

The layered geometry hypothesis reinterprets the evolution of the sound horizon and its imprint on baryon acoustic oscillations (BAO) within a unified geometric framework. Unlike the standard $\Lambda$CDM model, where the sound horizon's evolution is influenced by radiation density and early-universe expansion, this hypothesis situates the process within the exponentially expanding 3-sphere of the noumenological universe. The sound horizon forms within photon-baryon acoustic oscillations governed by local dynamics but reflects the global coherence of the noumenological structure. This layered interplay between local and global phenomena emphasizes the need for recalibrated density parameters that account for the transformation from noumenological to phenomenological perspectives.

Recalibrating $\Omega_r$, $\Omega_m$, and $\Omega_\Lambda$ ensures the model aligns with observed phenomena while capturing the distinct causal and geometric structures of the layered framework. Transformations between the noumenological universe's intrinsic expansion and the flat $\Lambda$CDM phenomenological expansion necessitate adjustments to the sound horizon's mapping onto BAO scales. These recalibrations are critical for resolving observed tensions in cosmological data, including the Hubble constant discrepancy and evolving dark energy signals, while avoiding reliance on ad hoc mechanisms.

High-precision BAO datasets, such as those from DESI \cite{Aghamousa2016} and Euclid \cite{Laureijs2011}, offer robust benchmarks for testing these predictions. By comparing recalibrated sound horizon features to observed transverse comoving distances and Hubble rates across redshifts, researchers can evaluate the hypothesis' capacity to resolve discrepancies inherent in $\Lambda$CDM. This analysis bridges early- and late-universe dynamics, linking recalibrated BAO features to the layered geometry's predictions, providing a decisive test for the framework's validity.

The hypothesis also offers a fresh perspective on dark energy by decoupling the noumenological and phenomenological descriptions of cosmic expansion. Unlike the standard $\Lambda$CDM model, where dark energy is treated as a uniform energy density driving accelerated expansion, the layered geometry hypothesis attributes this acceleration to the global properties of the noumenological 3-sphere. The flat $\Lambda$CDM expansion rate observed phenomenologically in the matter frame reflects the transformation of noumenological processes, rooted in the universe's intrinsic geometry.

This reinterpretation eliminates the need for a mysterious dark energy component while maintaining consistency with observational evidence, including CMB anisotropies and BAO. By leveraging datasets from Planck \cite{Aghanim2020}, DESI \cite{Aghamousa2016}, and Euclid \cite{Laureijs2011}, researchers can explore whether this alternative explanation for cosmic acceleration aligns with observed data.

The layered geometry hypothesis provides a compelling alternative to inflationary models for explaining key features of the cosmic microwave background (CMB). Traditional inflationary paradigms attribute the coherence, adiabatic nature, and near scale-invariance of the CMB anisotropies to a rapid period of exponential expansion driven by an inflaton field. By contrast, the layered framework derives these features directly from the noumenological exponential expansion of the 3-sphere, described by $\sim \cosh(\sqrt{\Lambda/3}T)$. This reinterpretation suggests that the observed properties of the CMB emerge naturally from the universe's intrinsic geometric structure, eliminating the need for an inflationary mechanism.

A notable feature of the CMB that challenges standard inflationary predictions is the suppression of power at low multipoles, particularly the quadrupole moment \cite{Komatsu2011, Aghanim2020}. The hypothesis attributes this suppression to constraints on initial perturbations imposed by the layered noumenological geometry. Similarly, the quadrupole-octupole alignment, a statistically significant deviation from isotropy \cite{Tegmark2003, Aghanim2020}, arises naturally from the interplay between the 3-sphere's topology and the phenomenological mapping to the matter frame.

Through rigorous theoretical modeling and comparison with high-precision observations, the hypothesis reframes the coherence, scale-invariance, and adiabatic nature of the CMB anisotropies as emergent properties of the layered geometry. By disentangling these features from inflationary assumptions, the model challenges the necessity of an inflaton field and offers a geometric explanation grounded in the universe's layered structure.

Validating these interpretations requires rigorous theoretical modeling and comparison with high-precision observations. Simulations of perturbation evolution under the hypothesis can predict the expected amplitude and alignment of these anomalies, offering a quantitative basis for testing. Observations from Planck \cite{Aghanim2020} and future missions like LiteBIRD \cite{Ghigna2024} provide the necessary precision to evaluate whether the layered framework reproduces these features more effectively than inflationary models.

Beyond anomalies, the hypothesis reframes the coherence, scale-invariance, and adiabatic nature of the CMB anisotropies as emergent properties of the layered geometry. These features, often viewed as evidence for inflation, are reinterpreted as consequences of the global coherence maintained by the noumenological exponential expansion. By disentangling these properties from inflationary assumptions, the model challenges the necessity of an inflaton field and offers a geometric explanation grounded in the universe's layered structure.

The reinterpretation of CMB anomalies underscores the hypothesis' potential to unify cosmological phenomena under a single geometric paradigm. Through detailed observational and theoretical analyses, this framework offers a pathway for addressing longstanding questions about the universe's early evolution while critically evaluating the role of inflationary mechanisms. The layered geometry hypothesis invites scrutiny as a cohesive and testable alternative, with the potential to reshape our understanding of the CMB and its origins.

The layered geometry hypothesis presents a robust and testable alternative to the standard FLRW paradigm, offering new pathways for understanding the universe's evolution. By reframing phenomena traditionally attributed to inflation, dark matter, and dark energy as emergent properties of a layered geometric structure, the model challenges entrenched assumptions while aligning with observational evidence. Its capacity to unify disparate observations under a single framework underscores its transformative potential, bridging gaps between theory and data.

Through its unique treatment of global and local dynamics, the hypothesis reinterprets key cosmological features---such as the coherence of the CMB anisotropies, the apparent gravitational effects of dark matter, and the accelerated expansion of the universe---as manifestations of a unified noumenological structure. These insights invite rigorous empirical scrutiny, offering a pathway to resolve persistent tensions in cosmology, such as the Hubble constant discrepancy, the sound horizon calibration problem, and CMB anomalies.

As upcoming datasets from DESI, Euclid, and LiteBIRD refine our understanding of cosmic structure, the layered framework provides a compelling avenue for addressing unresolved questions in modern cosmology. By integrating early- and late-universe phenomena within a single geometric paradigm, the hypothesis paves the way for a deeper understanding of the universe's fundamental nature. It challenges the community to critically reevaluate longstanding assumptions, presenting a cohesive and testable vision for the future of cosmological research.

\section{Discussion and Conclusions}\label{sec:conclusions}

This work addresses foundational tensions in physics by reexamining the concept of physical existence and reinterpreting its role within relativity and cosmology. By critically examining the linguistic, philosophical, and theoretical foundations of modern physics, this paper proposes a cosmological interpretation of relativity that bridges empirical evidence with conceptual clarity. This framework challenges entrenched assumptions and provides a coherent pathway to resolve long-standing problems in cosmology and fundamental physics.

A central insight of this work is the distinction between existing and happening, enduring and occurring, which addresses fundamental misconceptions about what it means for something to physically exist. While the three-dimensional universe exists and endures in the present, described mathematically through a four-dimensional space-time framework, space-time itself does not exist in the same sense. This distinction has been obscured by centuries of linguistic and philosophical imprecision, perpetuating a tacit conflation of phenomena with noumena. Misinterpretations of existence have introduced conceptual confusion, such as the implicit assumption of a non-physical meta-time absent from the theory of relativity. By disentangling these conflations, this work constructs a coherent framework that reconciles intuition with empirical evidence, clarifying the physical meaning of existence and its role in modern physics.

The paper highlights a critical dissonance in modern physics: the operational interpretation of simultaneity in Einstein's relativity is fundamentally at odds with empirical evidence from cosmology, which strongly supports a preferred cosmic rest frame and absolute cosmic time. While Einstein's operationalist framework was revolutionary for its time, prioritizing measurable quantities over ontological clarity, its limitations are increasingly evident in light of modern observations. This work demonstrates that Einstein's rejection of absolute simultaneity in favor of frame-dependent time conflicts with the global coherence evident in the large-scale structure of the universe.

By reinterpreting relativity through a cosmological lens, this work restores coherence between relativity and cosmology, redefining time as an objective physical quantity tied to the dynamical evolution of a three-dimensional universe. This redefinition challenges the foundational assumptions of relativistic physics, showing that the absence of absolute simultaneity in Einstein's interpretation introduces inconsistencies when applied to cosmological phenomena. Rejecting this operational definition reveals a more unified framework that aligns relativity with modern cosmological evidence, providing clarity and coherence to the relationship between time, space, and existence.

Gravitational collapse, a cornerstone of black hole physics, is fundamentally reinterpreted in this work, exposing a critical dependence on flawed simultaneity assumptions. The standard model's depiction of fully-formed black holes with event horizons and singularities, championed by Penrose's singularity theorems, rests on the implicit acceptance of synchronous slicing within space-time. This assumption, however, is neither required by general relativity nor consistent with observational evidence. General relativity leaves simultaneity between events on separate worldlines undefined, and interpreting them as such misrepresents the mathematical theory. This paper demonstrates that gravitational collapse is an inherently asymptotic process, with event horizons forming only at the end of infinite cosmic time.

This reinterpretation resolves paradoxes associated with singularities, reframing black holes as dynamical, evolving structures deeply integrated within the universe's layered geometry. Beyond rethinking individual black holes, this insight reveals a profound connection between gravitational collapse and cosmic origins. By embedding collapse within the layered framework, the model suggests a seamless transition between the end states of collapsing structures and the genesis of new universes. This novel perspective unifies local relativistic phenomena with global cosmological evolution, bridging the apparent divide between black hole physics and the broader dynamics of cosmic structure.

Building on these insights, this work develops a cosmological model grounded in a layered geometrical framework that bridges local dynamics and global coherence. By discarding the synchronicity assumption while preserving isotropy, homogeneity, and absolute cosmic time, the model avoids the conflation of phenomena with the universe's underlying structure. This reinterpretation provides a more fundamental perspective on cosmic evolution, revealing that the universe's intrinsic geometry underpins its observed properties without invoking inflationary mechanisms or fine-tuned initial conditions.

The layered geometry resolves foundational cosmological problems, including the flatness, horizon, and cosmological constant problems, as natural consequences of its framework. The observed flat $\Lambda$CDM expansion rate, traditionally linked to energy densities in FLRW cosmology, emerges phenomenologically from the universe's layered structure, independent of material content. By situating isotropy and homogeneity as emergent properties rather than built-in assumptions, the model realigns cosmology with relativistic principles, offering a coherent and geometrically grounded alternative to the standard paradigm.

A critical insight of this work is the recognition that cosmic expansion operates across interconnected layers, reconciling local gravitational dynamics with global coherence. This layered geometrical framework explains how local deviations from Hubble's law, driven by gravitational interactions and density variations, are nested within the globally invariant expansion described by the Hubble flow. A gravitational threshold governs this relationship: structures exceeding this threshold, such as galaxies and clusters, remain gravitationally bound, resisting expansion, while underdense regions below this threshold uniformly participate in cosmic expansion. At supercluster scales, as demonstrated in \citeauthor{Tully2014}'s~\citeyearpar{Tully2014} analysis of the Laniakea Supercluster (and extended through the expanded scale of multiple basins of attraction \cite{Valade2024}), this interplay becomes especially evident. Superclusters exhibit complex gravitational flows layered atop the uniform expansion of underdense regions, revealing a nuanced dynamic that the homogeneity assumption of FLRW cosmology cannot accommodate.

This layered perspective extends to the interplay between relativistic space-time and cosmic time. While local gravitational interactions warp space-time and the apparent passage of time, the layered geometry situates these effects within a broader cosmological framework. Cosmic time and space are described globally, capturing the interplay between local deviations and uniform expansion. By disentangling these layers, the model provides a cohesive framework that bridges the observed dynamics of local and global phenomena. This reinterpretation underscores the inadequacy of FLRW cosmology's foundational assumptions while offering a coherent alternative that respects relativistic principles and aligns with empirical observations.

The FLRW cosmological model, while foundational to modern physics, is revealed in this work to suffer from epistemological flaws akin to those of geocentrism. Central to these flaws is the synchronicity assumption, which conflates phenomena---observable properties like isotropy and homogeneity---with noumena, the universe's underlying structure. This assumption violates the relativistic principle of frame-dependence, embedding a simplistic interpretation of reality that obscures the layered dynamics of cosmic evolution. Consequently, FLRW's reliance on inflation to explain isotropy, homogeneity, and the flatness of space-time mirrors the ad hoc mechanisms of epicycles and equants in geocentrism: effective for fitting data but ultimately masking deeper foundational issues.

The layered geometry hypothesis resolves these tensions by treating phenomena as emergent properties of an underlying noumenological structure, fundamentally altering our understanding of cosmic evolution. In this model, isotropy and homogeneity emerge naturally from the interplay between local deviations and global coherence, while the flat $\Lambda$CDM expansion rate arises as a phenomenological consequence of the universe's intrinsic geometric structure. This framework eliminates the need for inflation or fine-tuned initial conditions, reinterpreting the cosmological constant, $\Lambda$, as a geometric property rather than an energetic anomaly. By shifting focus from ad hoc mechanisms to the intrinsic geometry of the universe, the layered geometry hypothesis offers a coherent and empirically grounded alternative to FLRW cosmology.

Furthermore, the layered framework aligns with the relativistic insights that frame-dependent phenomena often deviate from the underlying reality. By grounding cosmology in a geometrically coherent model that respects these insights, this work challenges the assumptions of FLRW and positions the layered geometry hypothesis as a more fundamental and empirically consistent paradigm.

Empirical validation of this framework offers an exciting opportunity to rigorously test fundamental assumptions in cosmology. The predictions of this model---from the recalibration of density parameters to the layered geometrical interpretation of the sound horizon and BAO scale---provide clear, falsifiable avenues for investigation. By applying this framework to existing datasets and adopting statistical measures to assess validity, researchers can critically explore its explanatory power alongside $\Lambda$CDM and related models. Such recalibration could illuminate whether current tensions, such as the Hubble constant discrepancy and indications of evolving dark energy, arise from deeper assumptions about early-universe dynamics.

Speculative but promising connections to particle physics further extend the model's implications. The dual causal structures of the noumenological geometry suggest potential explanations for matter-antimatter duality and asymmetry and the $U(1)$, $SU(2)$, and $SU(3)$ symmetries of particle physics. These connections hint at a unified framework where fundamental symmetries emerge from the cosmological geometry itself, providing a novel perspective on the intersection of particle physics and cosmology.

This work challenges the entrenched operational epistemology that dominates modern physics, demonstrating its limitations in addressing foundational tensions. Modern methodology often prioritizes hypothesis-driven empiricism, where ad hoc mechanisms are introduced to reconcile observed phenomena with entrenched paradigms. While this approach has delivered predictive utility, it has fostered a deep reluctance to interrogate the foundational assumptions of core theories. This inertia has walled physics off from exploring inconsistencies and perpetuated reliance on frameworks like FLRW cosmology, which conflate phenomena with noumena.

The reductive methodology employed here provides a powerful alternative, systematically disentangling phenomena from interpretive frameworks to expose critical flaws in operational assumptions. This approach revealed how Einstein's operationalism, while revolutionary, fails to align relativity with empirical cosmological evidence, necessitating a reinterpretation of simultaneity and cosmic time. Similarly, the critique of FLRW cosmology demonstrates that the operationalist mindset fosters reliance on ad hoc mechanisms, such as inflation and dark energy, obscuring the framework's foundational contradictions.

Most importantly, this work underscores the necessity of integrating reduction with empirical testing as complementary tools for scientific progress. Reduction goes beyond simply interpreting empirical data; it exposes hidden assumptions, clarifies their limitations, and redefines theories with greater coherence. Historically, this methodology has driven paradigm shifts, from Copernicus to Einstein, by challenging dogmatic assumptions and illuminating deeper truths. By embracing reduction alongside empiricism, physics can move beyond the flawed operationalist mindset which inherently conflates phenomena with noumena in spite of significant lessons to the contrary, and reinvigorate its quest for understanding, addressing the most profound questions about time, space, and the universe's nature.

This interdisciplinary approach integrates geometric, philosophical, and linguistic innovations, highlighting the importance of intellectual openness and conceptual clarity. By revisiting the logical structure of science, it demonstrates how reduction complements hypothesis testing as an essential tool for scientific progress. Just as past paradigm shifts---such as the Copernican revolution and Einstein's reconceptualisation of space-time---required the abandonment of entrenched dogmas, this work calls for a renewed commitment to reexamining physics' foundational principles. This commitment, rooted in rigorous reductive analysis, is essential for uncovering new physics and addressing the most profound questions about time, space, and physical laws.

This work represents the beginning of a broader call to action in physics: to rethink entrenched paradigms and embrace methodologies that prioritize understanding over convenience. Future research must rigorously test this model's predictions, refine its mathematical framework, and explore its potential to unify cosmology, relativity, and particle physics. The layered geometry hypothesis offers a pathway to address long-standing tensions, from the nature of time and existence to the interplay between local and global dynamics. By bridging theoretical elegance with empirical rigor, this framework exemplifies how reduction, coupled with interdisciplinary innovation, can illuminate new paths in physics. Embracing this approach may not only resolve foundational tensions but also redefine the principles by which we understand the universe, paving the way for a cohesive physical theory that unites geometry, epistemology, and ontology.

\section*{Acknowledgments}

I am grateful to Matt Strassler for an email discussion that helped frame the core arguments in this paper. I also thank Dan Reichart, Michael Fitzgerald, Kal\'{e}e Tock, Edward Gomez, Andrew Yen, David McKinnon, Saeed Salimpour, and Josh Haislip for valuable discussions and feedback on earlier drafts. ChatGPT was consulted extensively throughout the preparation of this paper, providing assistance in refining arguments, structuring sections, and enhancing clarity. This collaboration highlights the growing role of generative AI in research and writing. No additional financial support was received for the writing of this paper, and no data were produced.

\bibliography{refs}

\end{document}